\documentclass[twoside,11pt,nolof]{starlink}
% Used for PODs
\def\C++{{\rm C\kern-.05em\raise.3ex\hbox{\footnotesize ++}}}
\def\underscore{\leavevmode\kern.04em\vbox{\hrule width 0.4em height 0.3pt}}
%\setlength{\parindent}{0pt}

% +
%  Name:
%     sun269.tex

%  Purpose:
%     SUN documentation for ORAC-DR LCOGT pipeline data reduction (SUN/269)

%  Authors:
%     Tim Lister (Las Cumbres Observatory)

%  Copyright:
%     Copyright (C) 2014-2015 Las Cumbres Observatory Global Telescope Network Inc.
%     All Rights Reserved.

%  History:
%     $Log$


% -


% ? Specify used packages
%\usepackage{graphicx}        %  Use this one for final production.
% \usepackage[draft]{graphicx} %  Use this one for drafting.
% ? End of specify used packages

%\pagestyle{myheadings}

% -----------------------------------------------------------------------------
% ? Document identification
% Fixed part
\stardoccategory  {Starlink User Note}
\stardocinitials  {SUN}
\stardocsource    {sun\stardocnumber}
\stardoccopyright {Copyright \copyright\ 2014--2015 Las Cumbres Observatory Global Telescope Inc.}

% Variable part - replace [xxx] as appropriate.
\stardocnumber    {269.4}
\stardocauthors   {Tim Lister \\
		    Las Cumbres Observatory}
\stardocdate	   {February 2015}
\stardoctitle     {ORAC-DR -- LCOGT pipeline data reduction}
\stardocversion   {0.4}
\stardocmanual    {User Guide}
\stardocabstract  {ORAC-DR is a general-purpose automatic data-reduction
pipeline environment.  This document describes its use to reduce optical imaging
data collected on the Las Cumbres Observatory Global Telescope Network (LCOGT).
}

% ? End of document identification
% -----------------------------------------------------------------------------

% +
%  Name:
%     sun269.tex
%
%  Purpose:
%     Documentation for LCOGT data reduction with ORAC-DR
%
%  Authors:
%     TAL: Tim Lister (LCOGT)
%
%  History:
%     03-NOV-2014 (TAL):
%        Initial version circulated to check form and function.
%     18-JAN-2015 (TAL):
%        Updated version.
%     09-FEB-2015 (TAL):
%        Add appendices with index.dqc and index.astromqc details.
%     {Add further history here}
%
% -

\stardocname{\stardocinitials /\stardocnumber}
%\markboth{\stardocname}{\stardocname}
%\setlength{\textwidth}{160mm}
%\setlength{\textheight}{230mm}
%\setlength{\topmargin}{-2mm}
%\setlength{\oddsidemargin}{0mm}
%\setlength{\evensidemargin}{0mm}
%\setlength{\parindent}{0mm}
%\setlength{\parskip}{\medskipamount}
%\setlength{\unitlength}{1mm}

% ? Document specific \newcommand or \newenvironment commands.

\providecommand{\CCDPACK}{\xref{\textsc{{Ccdpack}}}{sun139}{}}
\providecommand{\GAIA}{\xref{\textsc{{Gaia}}}{sun214}{}}
\providecommand{\FIGARO}{\xref{\textsc{{Figaro}}}{sun86}{}}
\providecommand{\KAPPA}{\xref{\textsc{{Kappa}}}{sun95}{}}
\providecommand{\SMURF}{\xref{\textsc{{Smurf}}}{sun258}{}}
\providecommand{\SLALIB}{\xref{\textsc{{SLALIB}}}{sun67}{}}
\providecommand{\AASTROM}{\xref{\textsc{{AutoAstrom}}}{sun242}{}}
\providecommand{\ASTROM}{\xref{\textsc{{Astrom}}}{sun5}{}}
\providecommand{\ORACDR}{{\footnotesize ORAC-DR}}
\providecommand{\SEx}{\textsc{SExtractor}}
\providecommand{\SCAMP}{\textsc{SCAMP}}
\providecommand{\astref}{\xref{SUN/211}{sun211}{}}
\providecommand{\ndfref}{\xref{SUN/33}{sun33}{}}

% Telescopes and instruments
\providecommand{\AAT}{\htmladdnormallink{AAT}{http://www.aao.gov.au/}}
\providecommand{\JCMT}{\htmladdnormallink{JCMT}{http://www.jach.hawaii.edu/JACpublic/JCMT/}}
\providecommand{\UKIRT}{\htmladdnormallink{UKIRT}{http://www.jach.hawaii.edu/JACpublic/UKIRT/}}
\providecommand{\VLT}{\htmladdnormallink{VLT}{http://www.eso.org/instruments/}}


\providecommand{\oracdr}{\textsc{orac-dr}}
\providecommand{\oracsun}{\xref{SUN/230}{sun230}{}}
\providecommand{\scubasun}{\xref{SUN/231}{sun231}{}}
\providecommand{\picard}{\textsc{picard}}
\providecommand{\picardsun}{\xref{SUN/265}{sun265}{}}

% Application tasks
\providecommand{\task}[1]{\textsf{#1}}
\providecommand{\makebias}{\xref{\task{makebias}}{sun139}{MAKEBIAS}}
\providecommand{\makecal}{\xref{\task{makecal}}{sun139}{MAKECAL}}
\providecommand{\makeflat}{\xref{\task{makeflat}}{sun139}{MAKEFLAT}}
\providecommand{\debias}{\xref{\task{debias}}{sun139}{DEBIAS}}
\providecommand{\calcor}{\xref{\task{calcor}}{sun139}{CALCOR}}
\providecommand{\flatcor}{\xref{\task{flatcor}}{sun139}{FLATCOR}}

% macros for typesetting parameters
\providecommand{\aparam}[1]{\texttt{#1}}     % ADAM parameter
\providecommand{\cparam}[1]{\texttt{#1}}     % CONFIG parameter
\providecommand{\ndfcomp}[1]{\texttt{#1}}    % NDF component

% LCOGT things
\providecommand{\LCO}{LCOGT}
\providecommand{\LCONet}{LCOGT Network}
\providecommand{\LCOFull}{Las Cumbres Observatory Global Telescope Network}
\providecommand{\SAICD}{LCOGT-SA-ICD Version 0.10.0}
\providecommand{\SAICDFull}{LCOGT Pipeline/Science Archive Interface Control Document (LCOGT-SA-ICD Version 0.10.0)}
\providecommand{\SA}{Science Archive}
\providecommand{\Ast}{\texttt{timecorrect.py} code}


\providecommand{\rmn}[1] {\mathrm{#1}}
\providecommand{\arcmin}{\mbox{$^{\prime}$}}
\providecommand{\arcsec}{\mbox{$^{\prime \prime}$}}

% ? End of document specific commands
% -----------------------------------------------------------------------------
%  Title Page.
%  ===========
%\renewcommand{\thepage}{\roman{page}}
\begin{document}
%\setcounter{secnumdepth}{5}
%\thispagestyle{empty}

% ? Add picture here if required for the LaTeX version.
%   e.g. \includegraphics[scale=0.3]{filename.ps}
\startitlepic {\begin{center}
\begin{tabular}{cc}
\includegraphics[width=1.0in]{sun260_logo} & 
\includegraphics[height=1.0in]{sun269_logo} \\
\end{tabular}
\end{center}}
% ? End of picture
\scfrontmatter


%\cleardoublepage
%\renewcommand{\thepage}{\arabic{page}}
%\setcounter{page}{1}



\section{\xlabel{introduction}Introduction\label{sec:intro}}

\ORACDR\ is a data-reduction pipeline operating at \UKIRT, \JCMT, the \AAT\  and
forms the basis of the operational pipeline at LCOGT. The pipeline reduces and
displays multi-frame observations soon after they are read from the detector. 
This allows observers to assess the quality and suitability of their data in
near real time but \ORACDR\ is also used to produce the final archive-delivered
products using the same code base.

\xref{SUN/230}{sun230}{} presents an overview of \ORACDR,
general facilities like its display system, and it explains the
differences between a pipeline and a traditional reduction package.
Put briefly, \ORACDR\ uses a few data headers to direct the data
reduction.  Amongst these headers is the name of a \emph{recipe}.  A
recipe is a series of high-level instructions such as ``make a
mosaic'' or ``divide by a flat'' that reduces an \emph{observation\/}
comprising one or more data frames.  The implementation of each of
these instructions is through a Perl script---called a
primitive---which calls Starlink packages such as
\xref{\CCDPACK}{sun139}{} and \xref{\KAPPA}{sun95}{},
to actually do the processing of the bulk data.

This document describes how to use \ORACDR\ software on Starlink to reduce
optical imaging data from the \LCOFull\ (hereafter LCOGT; \cite{brown2013lcogt})
consists of data taken from one of the LCOGT instruments; \textsc{LCOSBIG},
\textsc{LCOCC}, \textsc{LCOSINISTRO}, \textsc{LCOSPECTRAL} or \textsc{LCOFLI}.

\subsection{Document conventions}

In an attempt to make this document clearer to read, different fonts
are used for specific structures.

Observing modes are denoted by all upper case body text (e.g.\
FLATFIELD).

Starlink package names are shown in small caps (e.g.\ \SMURF);
individual task names are shown in sans-serif
(e.g.\ \makebias). \oracdr\ recipes and primitives are also shown in
sans-serif and are always upper case (e.g.\ \task{REDUCE\_BIAS}).

Content listings are shown in fixed-width type (sometimes called
`typewriter'). Extensions and components within NDF (\ndfref) data
files are shown in upper case fixed-width type (e.g.\
\ndfcomp{HISTORY}).

Text relating to filenames (including suffices for data products), key
presses or entries typed at the command line are also denoted by
fixed-width type (e.g.\ \texttt{\% kappa}), as are parameters for
tasks which are displayed in upper case (e.g.\ \aparam{METHOD}).

References to Starlink documents, i.e., Starlink User Notes (SUN),
Starlink General documents (SG) and Starlink Cookbooks (SC), are given
in the text using the document type and the corresponding number
(e.g.\ SUN/95). Non-Starlink documents are cited in the text and
listed in the bibliography.

File name suffices represent the text between the final underscore
character and the three-letter \verb+.sdf+ extension. For example, a
file named \verb+lsc1m005-fl07-20140717-0173-d02_bp.sdf+ has the suffix
\verb+_bp+.

\section{\xlabel{pipelines}LCOGT Pipeline Variants\label{sec:pipelines}}

There are two variants of the LCOGT pipeline, one of which is
designed to run in real time as data comes in. Most users will only need to
run the second form, known as the offline pipeline.

\begin{itemize}
\item The quick-look (QL) pipeline is primarily designed to perform
  quality-assurance analysis of the incoming data for real-time
  assessment of the instrument performance. 

\item The offline pipeline has access to all the data observed for a
  given night/camera and adopts a best-possible reduction approach. Images
  are made for each complete observation which are combined to create
  the final image.
\end{itemize}

\subsection{Requirements for running the LCOGT pipeline}

The LCOGT pipeline requires a recent Starlink installation. The
latest version may be obtained from
\htmladdnormallink{\texttt{http://starlink.jach.hawaii.edu/starlink}}{http://starlink.jach.hawaii.edu/starlink}. Since
development of the pipeline is an ongoing process, it is recommended
that the newest builds be used to access the full capabilities of the
pipeline. These builds can be obtained from\\
\htmladdnormallink{\texttt{https://github.com/Starlink/ORAC-DR/}}{https://github.com/Starlink/ORAC-DR/}
and may be kept up-to-date with git.

The Starlink Perl installation (Starperl) must be used to run the
pipeline due to the module requirements. The Starlink environment should be
initialized as usual before running \oracdr.

The pipeline uses the following Starlink applications:
\begin{itemize}
\item \KAPPA
%\item \FIGARO
\item \CCDPACK
\item \SEx
\end{itemize}

\subsection{\xlabel{envvars}Important environment variables\label{sec:envvars}}

The pipeline uses a number of environment variables to determine where
data should be read from and written to. Some are set automatically
when the pipeline is initialized, but they can be overridden manually
and, with the \verb+-honour+ flag may be left unchanged between
successive runs of the pipeline. The variables that must be defined in
order for the pipeline to run are denoted as `Mandatory' in the list
below.

\begin{itemize}

\item \verb+STARLINK_DIR+: location of the user's Starlink
  installation. [Mandatory]

\item \verb+ORAC_DATA_IN+: the location where data are read from. If
  running with \verb+-loop flag+, this is the location of the flag
  files, rather than the data files. [Mandatory]

\item \verb+ORAC_DATA_OUT+: location where intermediate pipeline data products are
  written. [Mandatory]

\item \verb+FINAL_DATA_OUT+: final location for pipeline data products. At SBA,
  this is \\ \verb+/mfs-sba/engineering/<site>/<camera>+. If not defined, the current
  directory is assumed. [Optional]

\item \verb+ORAC_DATA_CAL+: a user-specified location for
  the library calibration files. [Mandatory]

\item \verb+ORAC_CAL_ROOT+: a user-specified location for
  pipeline configuration files. [Mandatory]

\item \verb+ORAC_INSTRUMENT+: the type of LCOGT instrument class in use e.g.
\task{LCOSINISTRO}. [Mandatory]

\item \verb+ORAC_CAMERACODE+: the type of LCOGT instrument instance in use e.g.
\task{fl07}. [Mandatory]

\item \verb+ORAC_SITECODE+: the LCOGT site code that the data is from  e.g.
\task{LSC}. [Mandatory]

\end{itemize}

\section{Running ORAC-DR}

This is a very brief introduction to running \ORACDR. More detailed
information can be found in \xref{SUN/230}{sun230}{}.
\xref{SUN/232}{sun232}{} also includes a description of how to set up
and run \ORACDR.

These instructions assume that you are using the \texttt{tcsh} shell as this is
what the pipeline wrappers have been written in and have the Starlink software
already setup (this is normally done in the users' setup files e.g. in
\texttt{~/.tcshrc}) but if this has not be done, you can get things setup by
issuing the following commands (this assumes the Starlink software has been
installed in \texttt{/git-star}):
\begin{terminalv}
      $ tcsh
      % setenv STARLINK_DIR /git-star
      % source $STARLINK_DIR/etc/cshrc
      % source $STARLINK_DIR/etc/login
\end{terminalv}

You must first initialise \ORACDR\ using
\texttt{$\sim$eng/Setup\_Night\_ORAC.csh <UTdate> <camera>}. This will prepare
\ORACDR\ to reduce data taken that night. If you wish to reduce a previous
nights data then you should specify the UT date on the command line, e.g.\
\texttt{oracdr\_lcosbig 20140506 kb74}. If necessary, you should set the
\texttt{\$ORAC\_DATA\_IN} and \texttt{ \$ORAC\_DATA\_OUT} environment variables
to the names of the directories from which the raw data should be read and to
which reduced data should be written.

For example:

\begin{terminalv}
      % source ~eng/Setup_Night_ORAC.csh 20140511 kb74
\end{terminalv}

This will setup the necessary environment variables (as detailed in
Section~\ref{sec:envvars}) to enable the pipeline to run properly. It is
also possible to redefine the enviroment variables to read and/or write from
different locations or to use a different calibration library for example. (To
see a list of all the enviroment variables that were setup, you could issue the
following command:
\begin{terminalv}
      % printenv | grep "ORAC\|FINAL"
      FINAL_DATA_OUT=/mnt/images/daydirs/kb74/20140511/oracproc
      ORAC_CAL_ROOT=/data/tlister/SVN/LCO_Pipeline/trunk/code/oracdr/cal
      ORAC_CAMERACODE=kb74
      ORAC_DATA_CAL=/data/tlister/SVN/LCO_Pipeline/trunk/code/oracdr/cal/lcosbig/kb74
      ORAC_DATA_IN=/mfs-sba/engineering/elp/kb74/20140511/raw
      ORAC_DATA_OUT=/speedy/eng/elp/proc/kb74/20140511
      ORAC_DIR=/data/tlister/SVN/LCO_Pipeline/trunk/code/oracdr/src/
      ORAC_INSTRUMENT=LCOSBIG
      ORAC_LOOP=flag -skip
      ORAC_NFS_OK=1
      ORAC_PERL5LIB=/data/tlister/SVN/LCO_Pipeline/trunk/code/oracdr/src/lib/perl5/
      ORAC_PERSON=tlister@lcogt.net
      ORAC_SITECODE=elp 
\end{terminalv}

To reduce all data for a given night you should run 

\begin{terminalv}
      oracdr -log f -verbose -file <filename of frames>
\end{terminalv}

This will not normally produce much (if any) output to the screen as the
pipeline is designed to be run non-interactively. The pipeline will reduce the
data in the order they are specified in the file given after the \verb+-file+
option, using the recipe name given in the image header. Details of the
processing will be written to a logfile called \texttt{.oracdr\_NNNN.log} where
\texttt{NNNN} is the current process ID. It is written to
\texttt{\$ORAC\_DATA\_OUT} and is a hidden file (use e.g. \verb+% ls -ltrA .oracdr*+ 
to see them with the most recent files at the bottom).

Additional information about the processing can be seen by starting
\verb+oracdr+ with the \verb+-log fx+ option which will open a X window for the
log which has the advantage that warnings and errors are written to different,
independently scrollable windows and the reduction can be paused and restarted.

%Several windows will (eventually) open: an \ORACDR\ text display, \GAIA\
%windows and \textsc{Kapview} windows (a collective term for various
%\KAPPA\ display tasks). The pipeline will reduce the data as they
%are stored to disk, using the recipe name in the image header.
%
%The pipeline is meant to run without interference from the observer.
%Thus, although you can use the various \GAIA\ tools to examine images,
%the pipeline should not need to be stopped and/or restarted. If,
%however, you do need to restart the pipeline then this can be done
%using the {\tt -from} option on the command line:
%
%\begin{terminalv}
%      oracdr -loop flag -from 19 -skip
%\end{terminalv}
%
%This will re-reduce frames from 19 onwards if they have previously
%been reduced, then continue to wait for new frames to arrive. The {\tt
%  -loop flag} tells it not to exit when it runs out of frames. When
%reducing data off-line this should be omitted. The {\tt -skip} tells it to
%skip missing observations.
%
%To re-reduce a group of
%previously stored frames you can use the {\tt -list} option to specify
%a list of frames separated by commas or ranges separated by colons:
%
%\begin{terminalv}
%      oracdr -list 15,18:20
%\end{terminalv}
%
You may choose to reduce your data with a recipe other than the one specified
in the file headers. If you wished to test data reduction with a new
\SCAMP-based astrometric solver you had written, you could specify the {\tt
OFFLINE\_REDUCTION\_SCAMP} recipe on the command line,
for example:

\begin{terminalv}
      oracdr -log fx -verbose -debug -file OFFLINE_REDUCTION_SCAMP
\end{terminalv}

(The \verb+-debug+ option will cause additional debug information such as each
primitive call and timing information to be written into 
\texttt{.oracdr\_NNNN.log} file and additionally into a \texttt{ORACDR.DEBUG}
file in \texttt{\$ORAC\_DATA\_OUT}).
%
%Further, advanced recipes are available to perform more advanced and rigorous
%data reduction than is normally carried out in real-time at the telescope. To
%use these recipes, supply the {\tt -recsuffix ADV} command-line option:
%
%\begin{terminalv}
%      oracdr -list 18:20 -recsuffix ADV
%\end{terminalv}

%The chief advantage to using the {\tt -recsuffix ADV} option instead of
%supplying the full recipe on the command-line is if the advanced recipe does
%not exist, the default standard recipe will be used instead. Thus you need not
%fear reducing pointing observations using a recipe designed for science
%targets.

To exit (or abort) \ORACDR\ click on `Exit' in the text log window, or
type {\tt [ctrl]-c} in the xterm. The command {\tt oracdr\_nuke} can
be used to kill all DR-related processes, should you be having serious
problems.


\section{An overview of the reduction}

The data reduction for LCOGT depends on the type of data being
reduced. Calibration observations (bias, dark, flat) are
reduced differently from science observations. 

\subsection{Bias observations}

Bias observations are used to remove the instrumental DC offset applied to the
CCD detector to ensure positive values at the ADC.

Masterbiases are created if there are enough ($>2$) raw bias frames in a
contiguous group (frames having the same `molecule number', \texttt{MOLNUM}) when
the final frame of the group is reached (i.e. when  \texttt{MOLFRNUM == 
FRMTOTAL}). The masterbias is created using the \task{REDUCE\_BIAS} recipe which
in turn calls the \task{\_MAKE\_BIAS\_FROM\_GROUP\_} primitive. This uses \CCDPACK's \makebias\ command using
the \aparam{method=median}\ combination method without zero offsetting
(\aparam{zero=false}) to combine the raw bias frames.

\subsection{Dark observations}

Dark observations are used to remove the thermal noise from the detector.

Masterdarks are created if there are enough ($>2$) raw dark frames in a
contiguous group (frames having the same `molecule number', \texttt{MOLNUM}) when
the final frame of the group is reached (i.e. when  \texttt{MOLFRNUM == 
FRMTOTAL}). The masterdark is created using the \task{REDUCE\_DARK} recipe which
in turn calls the \task{\_MAKE\_DARK\_FROM\_GROUP\_} primitive. This uses
\CCDPACK's \makecal\ command using the \aparam{method=median}\ combination
method with the \aparam{expose=} parameter set to a (created temporary) file of
the dark exposure time in order to scale the darks before combining the raw dark
frames.

\subsection{Flatfield observations}

Flat observations are used to remove the illumination variation across the field
of view and to correct for pixel to pixel sensitivity variations.

Masterflats are created if there are enough ($>3$) raw flat frames in a
contiguous group (frames having the same filter and `molecule number', \texttt{
MOLNUM}) when the final frame of the group is reached (i.e. when  \texttt{MOLFRNUM
==  FRMTOTAL}). The masterflat is created using the \task{SKY\_FLAT} recipe
which in turn calls the \task{\_MAKE\_FLAT\_FROM\_GROUP\_} primitive. This uses
\CCDPACK's \makeflat\ command to combine the raw flat frames. The combination
method and whether cleaning of bad pixels is desired are passed in from the \task{SKY\_FLAT} recipe in the call to the
\task{\_MAKE\_FLAT\_FROM\_GROUP\_} primitive and is used to set the
\aparam{method=}\ combination method and the \aparam{clean=}\ parameters, with the normal default being
\aparam{method=broadened} and \aparam{clean=f}\  (as it was found that cleaning
introduced clumps of bad pixels, particularly in the sharp edges between
Sinistro quadrants). The resulting master flat is normalised to unity.

As masterflats are constructed, the \task{\_FLAT\_QC\_} primitive will check to
see if it has found a matching pair of flats (masterflats having the same filter
and binning factors, taken from the same night (which is checked via \texttt{DAYOBS} in
the header)). Once a suitable pair is found, a difference frame is formed (using
\KAPPA's \task{sub}) and the sigma is computed. If the sigma is within
tolerance, the two flats are filed with calibration system for future use
(including making copies into the master calibration library location pointed to
by \texttt{\$ORAC\_DATA\_CAL}), otherwise the flats are removed from the
calibration system flat index.


\subsection{Imaging science observations}

This section describes the processing that happens for the science frames i.e.
the vast majority of the frames. Note that currently `standard' 
(\texttt{-s00}) frames are treated the same as science (\texttt{-e00}) frames.
The standard frames could (and there is some support within other ORAC-DR
instrument recipes) be treated differently if desired. This processing is
controlled by the \task{OFFLINE\_REDUCTION}\ recipe.

\subsubsection{Bias subtraction}

Bias subtraction and bad pixel masking are actually handled by the generic
\task{\_PREPARE\_SINGLE\_FRAME\_}\ primitive which is called from within the
recipe-specific \task{\_OFFLINE\_REDUCTION\_HELLO\_}\ primitive. The appropriate
bias and bad pixel mask to use are selected based on the normal nearest-in-time
suitable calibration frame that matches the frame's binning.

\subsubsection{Dark scaling and subtraction}

Dark scaling (the master dark has previously been scaled to an exposure time of
unity) to the science frames' exposure time and subtraction is performed by the
\task{\_SUBTRACT\_DARK\_NO\_THRESH\_}\ primitive. (The normal
\task{\_SUBTRACT\_DARK\_}\ primitve's additional thresholding behavior is not
appropriate for the LCOGT optical imaging data which is much less suspectible to
extra hot pixels). The appropriate
dark frame to use is selected based on the normal nearest-in-time
suitable calibration frame that matches the frame's binning.

\subsubsection{Flatfield normalization}

Flatfield division is performed by the \task{\_DIVIDE\_BY\_FLAT\_}\ primitive
and makes use of the \CCDPACK\ \flatcor\ task to perform the actual
correction and trimming to the trim limits. The appropriate
flatfield frame to use is selected based on the normal nearest-in-time
suitable calibration frame that matches the frame's binning and filter.

\subsubsection{Astrometric solution}

The current astrometric solver makes use of the \AASTROM\ wrapper around the
\ASTROM\ task and is controlled from within the \task{\_ADD\_AUTO\_ASTROMETRY\_}
primitive. This handles the detection and extraction of objects on the image
using \SEx\ and the download of a reference catalog to match against. The
reference catalogs are consulted in the order listed in the
\texttt{\$ORAC\_DATA\_CAL/index.skycat\_catalogue} and the default is UCAC-3
\cite{ucac3} with a fallback of 2MASS \cite{2mass} (both with two possible
server sources at CDS, France and CADC, Canada to add redundancy) and a final
fallback of the USNO-B \cite{usnob} catalog (The Context Cameras
(\texttt{LCOCC}) have a different version of
\texttt{\$ORAC\_DATA\_CAL/index.skycat\_catalogue} which only makes use of the
Tycho-2 catalog \cite{tycho2} due to the larger field of view and shallower
depth of these cameras).

The \ASTROM\ task can fit  4, 6, 7, 8, or 9 coefficients. The 4 coefficient is
never used (assumes same scale in both axes \& no shear terms), the 6
coefficient is the normal default for most cameras, the 7th term adds fitting
for the radial distortion (barrel/pincushion), the 8 coefficient fit adds
fitting for the plate center (but not distortion) and finally the 9 coefficient fit adds
fit includes fitting for the plate center along with the distortion and is the
default for the Context Cameras (\texttt{LCOCC}) due to their large field of
view.

After a putative fit, the results are sanity checked on the X and Y rms, the
number of stars used in the fit (must be $\geq$6) and the fitted platescale is
within 5\% of the expected value (which can cause false rejection if the nominal
platescale (\texttt{PIXSCALE} in the FITS header) placed by the Instrument Agent
hasn't been properly configured). The cases where the astrometric solution fails
or is reset after the sanity check are distinguished by setting \texttt{WCSERR}
in the header to '3' or '4' respectively (with '0' indicating success).

\subsubsection{Calculation of image statistics}

This occurs within the \task{\_CALCULATE\_SEEING\_STATS\_} primitive and makes
use of \SEx\ to perform the extraction and measurement of sources. A relatively
high threshold of 5.0 sigma above the background mesh fit with a minimum source
size of 12 pixels is used to avoid pushing down too far into the noise and
producing a biased result. Objects that have zero size or \SEx\
\texttt{FLAGS!=0} are removed from the calculation and then the mean FWHM,
ellipticity and orientation angle are calculated, reported in the log and filed
with the calibration system, appearing in \texttt{index.dqc}. These are subsequently
added to the header of the products as described in Section~\ref{sec:products}.

\subsubsection{Production of image catalog}

This occurs within the \task{\_CREATE\_IMAGE\_CATALOGUE\_}\ primitive and makes
use of \SEx\ to perform the source detection and extraction and produce a LCOGT
standard source catalog in the format as defined in the
\htmladdnormallink{\SAICDFull}{http://intranet.lco.gtn/Science_Archive\#Interface_Control_Document_.28ICD.29}. This
extraction occurs at a lower threshold of 1.5 sigma above the background mesh
fit to give a reasonably complete source catalog without overwhelming users or
the \SA\ with spurious detections.

The \SEx\ config file is in \texttt{\$ORAC\_DATA\_CAL/extractor\_catalogue.sex}
but some elements are dynamically modified within the primitive as the \SAICD\
specification calls for aperture photometry in fixed aperture of
1$^{\prime\prime}$,
3$^{\prime\prime}$, 5$^{\prime\prime}$ and 7$^{\prime\prime}$\ apertures which translates to a variable
camera-specific number of pixels, based on the platescale. The output catalog is
renamed to \texttt{$<$root name$>$10\_cat.fits} or \texttt{$<$root
name$>$90\_cat.fits} depending on whether QuickLook (indicated by
\texttt{\$Frm->uhdr( "QUICK\_LOOK" ) == 1}) or offline processing is occurring.
The output catalog is subsequently modified by the \Ast\ as further detailed in
Section~\ref{sec:timecorrect}.

\subsubsection{Calculation of zeropoint}

This makes use of the catalog of detected sources produced by the previous
\task{\_CREATE\_IMAGE\_CATALOGUE\_}\ primitive and performs a cross-match to the
combined standard star catalog produced by AJP's \task{findassm} program. The
cross match filters on distance, area (object size), ellipticity and quality
flags and surviving good matches are used to calculate the mean axial ratio of
bright sources,the sky brightness, the limiting magnitude, a transparency
estimate and the instrumental zero point. These results are reported in the log
and filed with the calibration system, appearing in \texttt{index.zeropoint}.
These are subsequently added to the header of the products as described in
Section~\ref{sec:products}.

\subsubsection{\xlabel{products}Output of derived products\label{sec:products}}

The \task{\_DERIVED\_PRODUCTS\_} primitive is used to create the derived products from the processed frames. The
following items are performed:
\begin{itemize}
\item  the headers are updated with the details of the master calibration frames
   and correction steps used,
\item  the saturation and trimmed section are updated with the values used in
   the pipeline processing,
\item  A clipped mean, median and sigma are calculated (using
\task{\_CLIPPED\_STATS\_MEDIAN\_}) and added to the header,
\item  the QC module (the \task{\_SET\_QC\_FLAGS\_}\ primitive) is run to produce the four QC flags in the headers,
\item  incorrect values of the TAGID and/or PROPID are flagged and corrected
   and the PROPID is checked against the lists of known proposals,
\item  access rights and state of the data are set,
\item  a 512$\times$512 PNG thumbnail of the processed image is created,
\item  the frame is converted to FITS format,
\item  the \Ast\ (see Section~\ref{sec:timecorrect}) is called to compute per-star barycentric time
   correction and airmass and updates the header of the BCD image and
   source catalog,
\item  if a final destination is defined by the enviroment variable
   \texttt{\$FINAL\_DATA\_OUT}, the derived products are moved there.
\end{itemize}

\subsubsection{\xlabel{timecorrect}Airmass and barycentric time
correction\label{sec:timecorrect}}


In order to provide the most precise absolute timestamp that is free from
ambiguities and limited by the properties of the target system, while allowing
for improved corrections in the future, we quote the site arrival time (and
timesystem) along with the correction to the Barycentric Julian Date in the
Barycentric Dynamical Time standard (${\rm BJD}_{\rm TDB}$; see \cite{eastman2010} for
further details). Although Barycentric Co-ordinate Time, TCB, is formally
recommended in the IAU framework (see \cite{iers2010}), it differs from the
formerly and widely used time systems such as TDB and TT by
$\sim0.5$\,s yr$^{-1}$ with a current difference of $\sim17$\,s which is likely to
lead to confusion for little benefit.

The per-star airmass and barycentric time correction is performed by
\task{timecorrect.py}\ which uses a Python-wrapped copy of the \SLALIB\ library
to perform the time and reference system transformations. This code calculates a
\ndfcomp{TCORR} column which is added to the source catalogs, before being
ingested into the archive. 

This \ndfcomp{TCORR} is the per-object complete ``time correction" to go from UTC start at the
observatory to TDB (Barycentric Dynamical Time) at the midpoint of the exposure
at the Solar System Barycenter (SSB). This \ndfcomp{TCORR} should be added to the UTC
start time to get a proper ${\rm BJD}_{\rm TDB}$. \ndfcomp{TCORR} is a collection of terms:
\begin{equation}
{\rm TCORR} = t_{\rm mid}+t_{\rm clock}+t_{\rm romer}+t_{\rm einstein}
\end{equation}
where $t_{mid}$ is half the exposure time (to get to the midpoint), $t_{clock}$
contains the $\rmn{UTC}\rightarrow\rmn{TDB}$ correction (which is time-dependent at the few ms level as
it involves knowing the Earth's orientation and the position of the observatory
within the Solar System potential), $t_{romer}$ is the R\o emer delay (light travel
time) from the observatory to the SSB and $t_{einstein}$ is the Einstein delay due
to the fact the observatory is moving (with the Earth) and displaced from the geocenter.

The aim of the code was that code should not contribute more than 1\% error on
the shortest exposure times likely on the \LCONet, which was taken to be 30\,Hz
($=33.3\,{\rm ms}$). The errors due to missing terms or limitations withing the
time correction code are believed to be at the few $\times\sim10\mu{\rm s}$
level due to:
\begin{itemize}
\item Missing from the above equation is the Shapiro delay (due to the light-bending
by the Sun's gravitational potential). This effect is of the order $\sim-20\,\mu\rmn{s}$
when observing an object 30 degrees from the Sun's limb, with the effect decreasing as the
Sun$\rightarrow$object angle increases.
\item The Earth's barycentric position is calculated using the \SLALIB\
\task{sla\_epv} routine which in turn uses a simplified version of the
\texttt{VSOP2000} planetary theory \cite{vsop2000}. Comparison by TAL with the
\texttt{INPOP10a} ephemeris \cite{inpop10a} over the range 2005--2031 produced
an RMS error of 3.83\ km or 12.8\ $\mu{\rm s}$.
\item The accuracy on the mean to apparent place transformation (used for
airmass) is sub-milliarcsecond, limited by the precession-nutation model (IAU
1976 precession, Shirai \& Fukushima 2001 forced nutation and precession
corrections \cite{sf2001}), which is also used to transform the Earth's position to
J2000.0 from the true equator and equinox of date.
\item The $\rmn{UTC}\rightarrow\rmn{TT}$ correction needs to know the correct
number of leapseconds which relies on the \SLALIB\ \task{sla\_dat} routine being
updated and the code recompiled when a new leap second is announced by the IERS.
\end{itemize}

The error on the time correction is most likely to be dominated by 1)
uncertainty on the NTP synchronization of the instrument control computer
($\sim$1--5\,ms), 2) uncertainty in the time-stamping of the exposure start
($\sim$100\,$\mu\rmn{s}$), and 3) uncertainty in the photon-weighted midtime of the
exposure.

\subsubsection{Photometric calibration}
Stefano's code...

\section{Data files}

\subsection{Filenames and locations}
Raw \LCO\ data conforms to the following format:\\
\texttt{<site><tel.\ class><tel.\ \#>-<inst>-YYYYMMDD-<nnnn>-<type><red.\ level>.fits}\\
where:
\begin{itemize}
\item{\texttt{<site>}} is the \LCO\ Network Site (one of \{ogg, coj, bpl, lsc, sqa, cpt, elp, tfn\}),
\item{\texttt{<tel.\ class>}} is the telescope class (one of \{2m0, 1m0, 0m8, 0m4\}),
\item{\texttt{<tel.\ \#>}} is the telescope number,
\item{\texttt{<inst>}} is the (unique) CCD descriptor,
\item{\texttt{YYYYMMDD}} is the UTC date at start of night (defined to be 16:00 local),
\item{\texttt{<nnnn>}} is the running number, restarting at 0001 when UTC date increments,
\item{\texttt{<type>}} is the type of file (from \{a=arc, b=bias, d=dark, e=exposure, f=sky flat, l=lambert flat, s=standard, x=experimental\}) 
\item{\texttt{<red.\ level>}} is the reduction level (00-99, currently one of
\{00 (raw data), 01 (flash processed), 02 (pre-processed), 10 (quicklook
processed) 90 (pipeline processed)\})
\end{itemize}

For ease of tracking data flow, the pipeline products will follow the same
policy, with the \texttt{<red.\ level>} altered appropriately to \texttt{90} in
the case of normally processed data. Catalogs derived from frames will inherit
the same rootname with a \texttt{\_cat.fits} suffix. Bitmap versions of the
frame will use the same rootname, replacing the \texttt{.fits} extension with
\texttt{.png}.

\subsection{File suffixes}

\textbf{Processing Frame suffixes}

\vspace{0.2cm}

\begin{tabular}{l c l}
\hline
Suffix & Kept & Description \hspace{9cm}  \\
\hline
\texttt{\_raw} & N & The raw frame\\
\texttt{\_bp}  & N & Bad pixels masked\\
\texttt{\_db}  & N & Bias frame removed\\
\texttt{\_bp\_db}  & N & Bad pixels masked and bias frame removed\\
\texttt{\_ff}  & N & Flat fielded\\
\hline
\end{tabular}

\vspace{0.5cm}

\textbf{Group suffixes}

\vspace{0.2cm}

Not used.

\vspace{0.5cm}

\textbf{FITS Processed Frame suffixes}

\vspace{0.2cm}

\begin{tabular}{l c l}
\hline
Suffix & Kept & Description \hspace{9cm}  \\
\hline
\texttt{\_\{a,b,d,e,f,g,w\}00} & Y & The raw frame\\
\texttt{\_\{a,b,d,e,f,g,w\}01}  & N & Pre-processed frame\\
\texttt{\_\{a,b,d,e,f,g,w\}02}  & N & Flash reduced\\
\texttt{\_\{e,s\}10}  & Y & QuickLook processed\\
\texttt{\_\{e,s\}90}  & Y & Offline processed\\
\hline
\end{tabular}


\begin{thebibliography}{}
\addcontentsline{toc}{section}{References}
\bibitem{brown2013lcogt} Brown, T.~M. et al., 2013, \textit{Las Cumbres
Observatory Global Telescope Network}, PASP, 125, 1031--1055

\bibitem{eastman2010} Eastman, J., Siverd, R. \& Gaudi, B. S., 2010, \textit{Achieving Better Than 1 Minute Accuracy in the Heliocentric and Barycentric
Julian Dates}, PASP, 122, 935

\bibitem{iers2010} Petit, G. \& Luzum, B. (eds), 2010, \textit{IERS Conventions
(2010)}, IERS Technical Note No. 36, Chap. 10

\bibitem{vsop2000} Moisson, X. \& Bretagnon, P., 2001, Celestial Mechanics and
Dynamical Astronomy, Volume 80, 3/4, 205--213

\bibitem{inpop10a} Fienga, A., Laskar, J., Kuchynka, P., Manche, H., Desvignes,
G., Gastineau, M., Cognard, I., Theureau, G., 2011, \textit{The INPOP10a
planetary ephemeris and its applications in fundamental physics}, 
Celestial Mechanics and Dynamical Astronomy, Volume 111, Issue 3, pp.363-385
DOI: 10.1007/s10569-011-9377-8

\bibitem{sf2001} Shirai, T. \& Fukushima, T., 2001, \textit{Construction of a
new forced nutation theory of the nonrigid Earth }, AJ, 121, 3270--3283. 

\bibitem{ucac3} Zacharias, N., Finch, C., Girard, T., Hambly, N.,  Wycoff, G.,
Zacharias, M.~I., Castillo, D., Corbin, T.,  DiVittorio, M., Dutta, S., Gaume,
R., Gauss, S.,  Germain, M., Hall, D., Hartkopf, W., Hsu, D.,  Holdenried, E.,
Makarov, V., Martinez, M., Mason, B.,  Monet, D., Rafferty, T., Rhodes, A.,
Siemers, T.,  Smith, D., Tilleman, T., Urban, S., Wieder, G.,  Winter, L.,
Young, A., 2010, \textit{The Third US Naval Observatory CCD Astrograph Catalog
(UCAC3)}, AJ, 139, 2184--2199.

\bibitem{2mass} Skrutskie, M.~F., Cutri, R.~M., Stiening, R.,  Weinberg, M.~D.,
Schneider, S., Carpenter, J.~M.,  Beichman, C., Capps, R., Chester, T., Elias,
J.,  Huchra, J., Liebert, J., Lonsdale, C., Monet, D.~G.,  Price, S., Seitzer,
P., Jarrett, T., Kirkpatrick, J.~D.,  Gizis, J.~E., Howard, E., Evans, T.,
Fowler, J.,  Fullmer, L., Hurt, R., Light, R., Kopan, E.~L.,  Marsh, K.~A.,
McCallon, H.~L., Tam, R., Van Dyk, S.,  Wheelock, S., \textit{The Two Micron All
Sky Survey (2MASS)}, 2006, AJ, 131, 1163--1183.

\bibitem{usnob} Monet, D.~G., Levine, S.~E., Canzian, B., Ables, H.~D., 
Bird, A.~R., Dahn, C.~C., Guetter, H.~H., Harris, H.~C.,  Henden, A.~A.,
Leggett, S.~K., Levison, H.~F.,  Luginbuhl, C.~B., Martini, J., Monet,
A.~K.~B.,  Munn, J.~A., Pier, J.~R., \textit{The USNO-B Catalog}, 2003, AJ,
125, 984--993.

\bibitem{tycho2} H{\o}g, E., Fabricius, C., Makarov, V.~V., Bastian, U., 
Schwekendiek, P., Wicenec, A., Urban, S., Corbin, T.,  Wycoff, G.,
\textit{Construction and verification of the Tycho-2 Catalogue}, 2000, A\&A,
357,367--386.

\end{thebibliography}

\newpage
\appendix

\section{\xlabel{ap_list}Alphabetical list of LCOGT recipes\label{ap:list}}

\begin{description}
\menuitem{OFFLINE\_REDUCTION}{
  Recipe for processing imaging science data}
\menuitem{QUICK\_LOOK}{
  QuickLook processing recipe for imaging science data}
\menuitem{REDUCE\_BIAS}{
  Process bias observations and form a masterbias}  
\menuitem{REDUCE\_DARK}{
  Process dark observations and form a masterdark}  
\menuitem{SKY\_FLAT}{
  Process sky flat observations and form a masterflat}
\end{description}

\newpage
\begin{small}
\section{\xlabel{ap_classified}Classified list of LCOGT primitives\label{ap:classified}}

LCOGT primitives may be classified in terms of their camera class and observing
mode as follows. Note that a lot of the \textbf{LCOSBIG} primitives are common to
the other instrument classes unless it has been overridden by an
instrument-specific version. The instrument-specific primitives are used in
preference to the generic ones when reducing that class of instrument data i.e.
the \textbf{LCOSINISTRO} version of \task{\_MAKE\_BIAS\_FROM\_GROUP\_} will be used
instead of the \textbf{LCOSBIG} one when reducing \textbf{LCOSINISTRO} data.


{\large
\begin{center}
\textbf{SBIG Camera (LCOSBIG) imaging data:}
\end{center}
}


{\large
\begin{center}
\textbf{Bias observations}
\end{center}
}
\begin{description}
%\classitem{\_AVERAGE\_BIASES\_}
%Averages the current frame into the current bias. (deprecated)
\classitem{\_FILE\_LIBRARY\_BIAS\_}
File a copy of the masterbias SDF, PNG and calibration index files in the
calibration library
\classitem{\_MAKE\_BIAS\_FROM\_GROUP\_}
Makes a masterbias from the current group of frames.
%\classitem{\_REDUCE\_BIAS\_}
%Reduce a group of bias frames for array tests to determine a group (deprecated)
\classitem{\_REDUCE\_BIAS\_HELLO\_}
Sets up data-reduction tasks and data for REDUCE\_BIAS recipes.
\classitem{\_REDUCE\_BIAS\_STEER\_}
Steers processing for REDUCE\_BIAS recipe.
\classitem{\_REDUCE\_BIAS\_TIDY\_}
Removes unwanted intermediate files for the REDUCE\_BIAS recipe.
\end{description}

{\large
\begin{center}
\textbf{Dark observations}
\end{center}
}
\begin{description}
\classitem{\_FILE\_LIBRARY\_DARK\_}
File a copy of the masterdark SDF, PNG and calibration index files in the
calibration library
\classitem{\_MAKE\_DARK\_FROM\_GROUP\_}
Makes a masterdark from the current group of frames.
\classitem{\_REDUCE\_DARK\_STEER\_}
Steers processing for REDUCE\_DARK recipe.
\classitem{\_REDUCE\_DARK\_TIDY\_}
Removes unwanted intermediate files for the REDUCE\_DARK recipe.
\classitem{\_REMOVE\_BIAS\_}
Subtracts a bias frame.
\end{description}

{\large
\begin{center}
\textbf{Flat observations}

\end{center}
}
\begin{description}
\classitem{\_FILE\_LIBRARY\_FLAT\_}
File a copy of the masterflat SDF, PNG and calibration index files in the
calibration library
\classitem{\_FLAT\_QC\_}
To perform Quality Control on flat fields produced by this pipeline.
\classitem{\_MAKE\_FLAT\_FROM\_GROUP\_}
Makes a masterflat from the current group of frames.
\classitem{\_REMOVE\_BIAS\_}
Subtracts a bias frame.
\classitem{\_SKY\_FLAT\_STEER\_}
Steers processing for SKY\_FLAT recipes.
\classitem{\_SUBTRACT\_DARK\_}
Subtracts a dark frame.
\classitem{\_SUBTRACT\_DARK\_NO\_THRESH\_}
Subtracts a dark frame without thresholding of remaining hot pixels.
\end{description}

{\large
\begin{center}
\textbf{Science observations}
\end{center}
}
\begin{description}
\classitem{\_ADD\_AUTO\_ASTROMETRY\_}
Performs automated astrometric corrections.
\classitem{\_CALCULATE\_IMAGE\_STATISTICS\_}
Calculate various image quality statistics based on an input catalog(ue).
\classitem{\_CALCULATE\_SEEING\_STATS\_}
Extract objects and determine the average FWHM, ellipticity, and position angle across a field.
\classitem{\_CALCULATE\_ZEROPOINT\_}
Calculate various image quality statistics based on an input catalogue.
\classitem{\_CREATE\_IMAGE\_CATALOGUE\_}
Performs source extraction and photometry on all sources.
\classitem{\_DATA\_QC\_TEST\_}
Evaluates the image data quality control flag and sets the relevant bitmask keywords in the data catalogue product.
\classitem{\_DERIVED\_PRODUCTS\_}
Create derived products from the processed frames.
\classitem{\_DERIVED\_PRODUCTS\_STEER\_}
Steers processing for DERIVED\_PRODUCTS recipe.
\classitem{\_DIVIDE\_BY\_FLAT\_}
Divides an object frame by a flatfield frame.
\classitem{\_IMG\_QC\_READOUT\_TEST\_}
Primitive description
\classitem{\_REMOVE\_BIAS\_}
Subtracts a bias frame.
\classitem{\_SUBTRACT\_DARK\_}
Subtracts a dark frame.
\classitem{\_SUBTRACT\_DARK\_NO\_THRESH\_}
Subtracts a dark frame without thresholding of remaining hot pixels.
\classitem{\_SET\_QC\_FLAGS\_}
Evaluate and set Quality Control flags.
\end{description}

{\large
\begin{center}
\textbf{Miscellaneous}
\end{center}
}
\begin{description}
\classitem{\_ADD\_PIPELINE\_VERSION\_}
Adds the current pipeline version to the file.
\classitem{\_CHECK\_PROPID\_}
Checks whether the proposal id in the frame header is known to the science archive.
\classitem{\_CLIPPED\_STATS\_MEDIAN\_}
Finds the clipped mean, median and standard deviation of a frame.
\classitem{\_CONVERT\_TO\_FITS\_}
Converts current observation to FITS.
\classitem{\_CREATE\_GRAPHIC\_FROM\_FILE\_}
Create a PNG, GIF, or JPG graphic from a given file.
\classitem{\_CREATE\_RAW\_FRAME\_}
Creates a raw frame in ORAC\_DATA\_OUT.
\classitem{\_DELETE\_RAW\_FRAME\_}
Remove the raw frame file.
\classitem{\_GET\_FILTER\_PARAMETERS\_}
Returns LCOGT optical filter characteristics.
\classitem{\_GET\_GAIN\_}
Finds the LCOGT instrument gain in electrons per ADU for the current Frame.
\classitem{\_GET\_READNOISE\_}
Finds the LCOGT instrument readnoise in electrons for the current Frame.
\classitem{\_GET\_SATURATION\_LEVEL\_}
Finds the LCOSBIG saturation level in ADU the current Frame.
\classitem{\_IMBOX\_STATS\_}
Primitive description
\classitem{\_INSTRUMENT\_HELLO\_}
Performs the LCOSBIG-specific imaging setup.
\classitem{\_LCOGT\_STANDARD\_MAGNITUDE\_}
Obtains the catalogue magnitude of a LCOGT standard.
\classitem{\_NIGHT\_LOG\_}
Produces a text file log of a night's imaging observations.
\classitem{\_OFFLINE\_REDUCTION\_HELLO\_}
Sets up data-reduction tasks and data for OFFLINE\_REDUCTION recipes.
\classitem{\_QUICK\_LOOK\_HELLO\_}
Sets up data-reduction tasks and data for QUICK\_LOOK recipes.
\classitem{\_SET\_ORIGIN\_}
Sets the origin of an observation.
\classitem{\_STANDARD\_MAGNITUDE\_}
Obtains the catalogue magnitude of a standard.
\classitem{\_TURN\_ON\_HISTORY\_}
Switches on history recording.
\end{description}

{\large
\begin{center}
\textbf{Context Camera (LCOCC) imaging data:}
\end{center}
}

\begin{description}
\classitem{\_CALCULATE\_ZEROPOINT\_NULL\_}
Null primitive for Context Camera
\classitem{\_GET\_FILTER\_PARAMETERS\_}
Return central wavelength, extinction and zeropoint for a given filter for the
ContextCamera.
\end{description}

{\large
\begin{center}
\textbf{Sinistro Camera (LCOSINISTRO) imaging data:}
\end{center}
}

\begin{description}
\classitem{\_CALCULATE\_ZEROPOINT\_}
alculate various image quality statistics based on an input catalogue.
\classitem{\_CREATE\_IMAGE\_CATALOGUE\_}
source extraction and photometry on all sources.
\classitem{\_DERIVED\_PRODUCTS\_}
Create derived products from the processed frames.
\classitem{\_FLAT\_QC\_}
To perform Quality Control on flat fields produced by this pipeline.
\classitem{\_GET\_FILTER\_PARAMETERS\_}
Returns LCOGT optical filter characteristics.
\classitem{\_GET\_GAIN\_}
Finds the LCOGT Sinistro instrument gain in electrons per ADU for the current Frame.
\classitem{\_GET\_READNOISE\_}
Finds the LCOGT Sinistro instrument readnoise in electrons for the current Frame.
\classitem{\_GET\_SATURATION\_LEVEL\_}
Finds the LCOSINISTRO saturation level in ADU the current Frame.
\classitem{\_MAKE\_BIAS\_FROM\_GROUP\_}
Makes a masterbias from the current group of frames.
\classitem{\_MAKE\_DARK\_FROM\_GROUP\_}
Makes a masterdark from the current group of frames.
\classitem{\_MAKE\_FLAT\_FROM\_GROUP\_}
Makes a masterflat from the current group of frames.
\end{description}

{\large
\begin{center}
\textbf{Spectral Instruments Camera (LCOSPECTRAL) imaging data:}
\end{center}
}

\begin{description}
\classitem{\_CREATE\_RAW\_FRAME\_}
Creates a raw frame in ORAC\_DATA\_OUT, removing legacy header items if present.
\classitem{\_GET\_GAIN\_}
Finds the LCOGT Spectral instrument gain in electrons per ADU for the current Frame.
\classitem{\_GET\_READNOISE\_}
Finds the LCOGT Spectral instrument readnoise in electrons for the current Frame.
\classitem{\_INSTRUMENT\_HELLO\_}
Performs the LCOSPECTRAL-specific imaging setup.
%\classitem{\_REDUCE\_BIAS\_}
%Reduce a group of bias frames for array tests to determine a group (deprecated)
\end{description}

{\large
\begin{center}
\textbf{FLI Autoguider Camera (LCOFLI) imaging data:}
\end{center}
}

\begin{description}
\classitem{\_ADD\_PIPELINE\_VERSION\_}
Adds the current pipeline version to the file.
\classitem{\_CREATE\_RAW\_FRAME\_}
Creates a raw frame in ORAC\_DATA\_OUT, removing legacy header items if present.
\classitem{\_MAKE\_DARK\_FROM\_GROUP\_}
Makes a masterdark from the current group of frames.
\end{description}

{\large
\begin{center}
{\it FLOYDS Andor Camera (LCOFLOYDS) spectroscopy data:}
\end{center}
}

\begin{description}
\classitem{\_ADD\_PIPELINE\_VERSION\_}
Adds the current pipeline version to the file.
\classitem{\_DETERMINE\_NREADS\_}
Determine the number of array reads per exposure.
\classitem{\_DETERMINE\_WAVEBAND\_}
Determine the waveband for flux calibration and blackbody calculations.
\classitem{\_FIX\_EXTRA\_HEADERS\_}
Null primitive for \textsc{LCOFLOYDS}.
\classitem{\_MAKE\_BIAS\_FROM\_GROUP\_}
Makes a masterbias from the current group of frames.
\classitem{\_MAKE\_RAW\_FILE\_}
Copies raw data to a \_mraw file.
%\classitem{\_REDUCE\_BIAS\_}
%Primitive description
\classitem{\_REDUCE\_BIAS\_HELLO\_}
Sets up data-reduction tasks and data for the \task{REDUCE\_BIAS} recipe.
\classitem{\_REDUCE\_BIAS\_STEER\_}
Steers processing for the \task{REDUCE\_BIAS} recipe.
\classitem{\_REDUCE\_BIAS\_TIDY\_}
Removes unwanted intermediate files for the \task{REDUCE\_BIAS} recipe.
\end{description}


\end{small}

\newpage

\section{\xlabel{ap_full_recipes}Specifications of LCOGT recipes\label{ap:full_recipes}}

The following pages describe the LCOGT recipes in detail. Recipes in
general live in subdirectories under \verb+$ORAC_DIR/recipes+ (this is the
general location for all recipes for imaging, spectroscopy, IFU's and general
purposes). For LCOGT use, the majority of recipes live under
\verb+$ORAC_DIR/recipes/imaging/LCOSBIG+ or a camera-class specific
subdirectory (e.g. \verb+$ORAC_DIR/recipes/imaging/LCOSINISTRO+). \oracdr\ 
searches for recipes via a `recipe search path' (controlled by the
\verb+orac_determine_recipe_search_path()+ subroutine in
\verb+ORAC::Inst::Defn+) which is ordered in
decreasing order of specificity e.g. for \textsc{LCOSINISTRO} the order will be:
\begin{itemize}
\item \verb+$ORAC_DIR/recipes/LCOSINISTRO/+
\item \verb+$ORAC_DIR/recipes/imaging/LCOSINISTRO/+
\item \verb+$ORAC_DIR/recipes/LCOSBIG/+
\item \verb+$ORAC_DIR/recipes/imaging/LCOSBIG/+
\item \verb+$ORAC_DIR/recipes/imaging/+
\end{itemize}

The recipe search path can also be overridden with an environment variable if
needed or a specific recipe to be run can be specified on the command line call
to \verb+oracdr+.

\newpage
\sstroutine{
   OFFLINE\_REDUCTION
}{
   Reduce a normal imaging EXPOSE observation
}{
   \sstdescription{
      It performs a debiassing, bad-pixel masking, dark subtraction,
      flat-field division, automatic astrometry, object detection and
      catalogue creation and converts to a PNG graphic and FITS version
      of the reduced frame. See the {\tt "}Notes{\tt "} for further information.
   }
   \sstnotes{
      All intermediate frames are deleted.
   }
}
\newpage
\sstroutine{
   QUICK\_LOOK
}{
   Perform a basic reduction of EXPOSE observations to give a {\tt "}quick
   look{\tt "} at the data
}{
   \sstdescription{
      It performs a debiassing, bad-pixel masking, dark subtraction,
      flat-field division, automatic astrometry, object detection and
      catalogue creation and converts to a PNG graphic and FITS version
      of the reduced frame. See the {\tt "}Notes{\tt "} for further information.
   }
   \sstnotes{
      All intermediate frames are deleted. The resulting output PNG
      graphic and FITS frames and catalog files have -\{e,s\}10.fits
      suffixes, rather than the normal -\{e,s\}90.fits
   }
}
\newpage
\sstroutine{
   REDUCE\_BIAS
}{
   Reduce a imaging BIAS observation
}{
   \sstdescription{
      Reduces a imaging BIAS observation, including coadding multiple
      integrations. Files the reduced bias frame for use by subsequent
      reduction of STARE and CHOP mode data.
   }
   \sstnotes{
      Creates a variance array for the bias frame, determined from the
      variance of the multiple integrations in the bias observation.
      There should be at least 3 integrations in a BIAS observation.
   }
   \sstdiytopic{
      See Also
   }{
      The \_REDUCE\_BIAS\_ imaging primitive.
   }
}
\newpage
\sstroutine{
   REDUCE\_DARK
}{
   Averages and files observations as the current dark
}{
   \sstdescription{
      This recipe reduces dark-frame observations from optical imaging
      data. Individual darks are scaled by the exposure time before being averaged.
      It files the single or averaged dark in the dark index file. Other
      reduction steps comprise bad-pixel masking, optional creation of
      data errors.
   }
   \sstnotes{
      \sstitemlist{

         \sstitem
         The bad-pixel mask applied is the binning-appropriate match taken from 
	 the \\ 
         \$ORAC\_DATA\_CAL/index.mask calibration index.

         \sstitem
         Intermediate frames are deleted.
      }
   }
   \sstdiylist{
      Configurable Steering Parameters
   }{
      \sstsubsection{
         USEVAR = LOGICAL
      }{
         Whether or not to create and propagate variance arrays. [0]
      }
   }
   \sstdiytopic{
      Output Data
   }{
      \sstitemlist{

         \sstitem
         The dark called dark\_$<$instrument$>$\_$<$UT
         night$>$\_$<$Xbin$>$\_$<$Ybin$>$ where $<$instrument$>$ is the LCOGT instrument identifier and $<$Xbin$>$,
         $<$Ybin$>$ are the binning factors in X and Y.

         \sstitem
         The dark is filed in \$ORAC\_DATA\_OUT/index.dark after checking whether it's
         mean is within the allowed tolerance.
         
         \sstitem
         If \$FINAL\_DATA\_OUT is set, the masterdark is converted to FITS and 
         copied there and a PNG version of the masterdark is created and moved 
         to \$FINAL\_DATA\_OUT.

      }
   }
   \sstimplementationstatus{
      \sstitemlist{

         \sstitem
         The processing engines are from the Starlink packages KAPPA and CCDPACK.

         \sstitem
         Uses the Starlink NDF format.

         \sstitem
         History is recorded within the data files.

         \sstitem
         The title of the data is propagated through the intermediate
         file to the dark.
      }
   }
}
\newpage
\sstroutine{
   SKY\_FLAT
}{
   Creates and files a flat field derived from dithered frames
}{
   \sstdescription{
      This recipe makes a sky flat for LCOGT imaging from a series of
      sky or Lambert flat frames which are combined using one of a
      selection of statistics. It expects a series of dithered sky
      frames in different filters.

      It performs debiassing, bad-pixel masking, and dark subtraction
      before combining normalised frames pixel by pixel using the
      median. Details of the flat are filed in the index of flats for
      future selection and use of the flat. See the \texttt{"}NOTES\texttt{"}  for further
      details.

      For best results the field observed should contain few stars and
      no bright ones. Ideally the telescope should have been pointed to
      the \texttt{"}flat spot\texttt{"}  - usually 5 to 10 degrees away from the zenith in
      the anti-solar direction to minimize gradients in the resultant
      flat.
   }
   \sstnotes{
      \sstitemlist{

         \sstitem
         The bad-pixel mask applied is the binning-appropriate match
         taken from the \\
         \$ORAC\_DATA\_CAL/index.mask calibration index.

         \sstitem
         Intermediate frames are deleted.
      }
   }
   \sstdiylist{
      Configurable Steering Parameters
   }{
      \sstsubsection{
         NUMBER = INTEGER
      }{
         The number of frames in the jitter. If absent, the number of
         offsets, as given by header NOFFSETS, minus one is used. If
         neither is available, 5 is used. An error state arises if the
         number of dithered frames is fewer than 3. []
      }
      \sstsubsection{
         USEVAR = LOGICAL
      }{
         Whether or not to create and propagate variance arrays. [0]
      }
   }
   \sstdiytopic{
      Output Data
   }{
      \sstitemlist{

         \sstitem
         The created flat fields are of the form \\
         flat\_$<$instrument$>$\_$<$UT
         night$>$\_$<$flattype$>$\_bin$<$Xbin$>$x$<$Ybin$>$\_$<$filter$>$[\_c$<$cycle\_number$>$]
         where $<$instrument$>$ is the LCOGT instrument identifier, $<$flattype$>$
         is the type of flat (either \texttt{'}SKYFLAT\texttt{'}  or \texttt{'}FAKEFLAT\texttt{'}  (a unity
         flat), $<$filter$>$ is the common name of the filter (e.g. \texttt{'}R\texttt{'} ) and
         $<$Xbin$>$, $<$Ybin$>$ are the binning factors in X and Y. If a flat of
         the same name already exists, a $<$cycle\_number$>$ will be added,
         counting from one.

         \sstitem
         The flats are filed in \$ORAC\_DATA\_OUT/index.flat.
      }
   }
   \sstimplementationstatus{
      \sstitemlist{

         \sstitem
         The processing engines are from the Starlink packages: CCDPACK,
         KAPPA, and FIGARO.

         \sstitem
         Uses the Starlink NDF format.

         \sstitem
         History is recorded within the data files.

         \sstitem
         The title of the data is propagated through the intermediate
         file to the dark.

         \sstitem
         Error propagation is controlled by the USEVAR parameter.
      }
   }
}

\newpage

\section{\xlabel{ap_full_primitives}Specifications of LCOGT primitives\label{ap:full_primitives}}

The following pages describe the LCOGT primitives in detail. Primitives in
general live in subdirectories under \verb+$ORAC_DIR/primitives+ (this is the
general location for all primitives for imaging, spectroscopy, IFU's and general
purposes). For LCOGT use, the majority of primitives live under
\verb+$ORAC_DIR/primitives/imaging/LCOSBIG+ or a camera-class specific
subdirectory (e.g. \verb+$ORAC_DIR/primitives/imaging/LCOSINISTRO+). \oracdr\ 
searches for primitives via a `primitive search path' (controlled by the
\verb+orac_determine_primitive_search_path()+ subroutine in
\verb+ORAC::Inst::Defn+) which is ordered in
decreasing order of specificity e.g. for \textsc{LCOSINISTRO} the order will be:
\begin{itemize}
\item \verb+$ORAC_DIR/primitives/LCOSINISTRO/+
\item \verb+$ORAC_DIR/primitives/imaging/LCOSINISTRO/+
\item \verb+$ORAC_DIR/primitives/LCOSBIG/+
\item \verb+$ORAC_DIR/primitives/imaging/LCOSBIG/+
\item \verb+$ORAC_DIR/primitives/imaging/+
\item \verb+$ORAC_DIR/primitives/general/+
\end{itemize}

The primitive search path can also be overridden with an environment variable if
needed.

\newpage
\sstroutine{
   \_ADD\_AUTO\_ASTROMETRY\_
}{
   Performs automated astrometric corrections
}{
   \sstdescription{
      This primitive automatically corrects astrometry for a given
      observation. It does so by downloading a catalogue from a source
      (typically UCAC-3), detects objects in the observation, correlates
      between the two catalogues, and calculates an astrometric
      solution.
   }
   \sstnotes{
      \sstitemlist{

         \sstitem
         This primitive is suitable for optical imaging instruments.

         \sstitem
         Processing only occurs when the steering header ADD\_ASTROMETRY
         is true.

         \sstitem
         Processing occurs on the current Group object.

         \sstitem
         Astrometric correction will probably fail if fewer than five
         objects are detected in the frame. Should this occur, the original
         WCS will be retained.

         \sstitem
         Should an astrometric solution be found, the WCS in the current
         Group will be overwritten with the solution.
      }
   }
   \sstdiytopic{
      Required Perl Modules
   }{
      Starlink::Autoastrom.
   }
}
\newpage
\sstroutine{
   \_ADD\_PIPELINE\_VERSION\_
}{
   Adds the current pipeline version to the file
}{
   \sstdescription{
      This primitive adds the current pipeline version to the file
   }
   \sstnotes{
      \sstitemlist{

         \sstitem
            This primitive is suitable for LCOGT optical imaging
            instruments.

         \sstitem
            This primitive needs to have SVN version properties set to
            update properly
      }
   }
   \sstdiytopic{
      Tasks
   }{
      KAPPA
         FITSMOD
   }
}
\newpage
\sstroutine{
   \_CALCULATE\_IMAGE\_STATISTICS\_
}{
   Calculate various image quality statistics based on an input
   catalogue
}{
   \sstdescription{
      This primitive calculates image quality statistics based on an
      input catalogue. This catalogue is typically the output from
      either Starlink::Extractor or Starlink::Autoastrom, so that object
      morphology information is available for calculations.

      This primitive calculates the mean axial ratio of bright sources,
      the average object diameter converted to K-band equivalent, the
      sky brightness, the limiting magnitude, and the instrumental zero
      point.
   }
   \sstnotes{
      \sstitemlist{

         \sstitem
         Files results with the Calibration system. FWHM is obtained
         with the fwhm() method, axial ratio is obtained with the
         axial\_ratio() method, sky brightness is obtained with the
         sky\_brightness() method, limiting magnitude is obtained with the
         limiting\_mag() method, and the zero point is obtained with the
         zeropoint() method.
      }
   }
   \sstdiytopic{
      Output Data
   }{
      None.
   }
   \sstdiytopic{
      Required Perl Modules
   }{
      Astro::Catalog.
   }
}
\newpage
\sstroutine{
   \_CALCULATE\_SEEING\_STATS\_
}{
   Extract objects and determine the average FWHM, ellipticity, and
   position angle across a field
}{
   \sstdescription{
      For the group file, this primitive finds good sources in the field
      and calculates the average FWHM, ellipticity, and position angle.
      It then displays these averages along with errors.
   }
   \sstnotes{
      \sstitemlist{

         \sstitem
         Currently uses SEXTRACTOR for source extraction.

         \sstitem
         The SEXTRACTOR configuration file is found in
         \$ORAC\_DATA\_CAL/extractor\_seeing\_stats.sex.

         \sstitem
         This primitive runs only when the CALCULATE\_SEEING\_STATS
         internal header is true.

         \sstitem
         This primitive only operates on the current Group file.
      }
   }
}
\newpage
\sstroutine{
   \_CALCULATE\_ZEROPOINT\_
}{
   Calculate various image quality statistics based on an input
   catalogue
}{
   \sstdescription{
      This primitive calculates image quality statistics based on an
      input catalogue. This catalogue is typically the output from
      either Starlink::Extractor or Starlink::Autoastrom, so that object
      morphology information is available for calculations.

      This primitive calculates the mean axial ratio of bright sources,
      the sky brightness, the limiting magnitude, a transparency estimate 
      and the instrumental zero point.
   }
   \sstnotes{
      \sstitemlist{

         \sstitem
         Files results with the Calibration system. FWHM is obtained
         with the fwhm() method, axial ratio is obtained with the
         axial\_ratio() method, sky brightness is obtained with the
         sky\_brightness() method, limiting magnitude is obtained with the
         limiting\_mag() method, and the zero point is obtained with the
         zeropoint() method.
      }
   }
   \sstdiytopic{
      Output Data
   }{
      None.
   }
   \sstdiytopic{
      Required Perl Modules
   }{
      Astro::Catalog.
   }
}
\newpage
\sstroutine{
   \_CHECK\_PROPID\_
}{
   Checks whether the proposal id in the frame header is known to the
      science archive
}{
   \sstdescription{
      The primitive checks whether the proposal id in the frame header
      is known to the IPAC Science Archive or the LCOGT Proposal DB,
      depending on the setting of \$test\_proposaldb.
   }
   \sstnotes{
      \sstitemlist{

         \sstitem
            This primitive is suitable for imaging cameras.

         \sstitem
            Processing only occurs for object frames.

         \sstitem
            Email sending requires Mail::Sendmail to be installed

         \sstitem
            The files to be checked are \$ORAC\_CAL\_ROOT/lcogt\_propid.txt or\\
            \$ORAC\_CAL\_ROOT/proposal\_list depending on whether the IPAC
            Science Archive or the LCOGT Proposal DB is being checked.
      }
   }
   \sstdiytopic{
      Required Perl Modules
   }{
      Mail::Sendmail
   }
}
\newpage
\sstroutine{
   \_CLIPPED\_STATS\_MEDIAN\_
}{
   Finds the clipped mean, median and standard deviation of a frame
}{
   \sstdescription{
      Use progressive sigma-clipping to find a representative mean,
      median and standard deviation of a frame. The default clipping
      thresholds give a reasonable approximation to the mode.
   }
   \sstnotes{
      \sstitemlist{

         \sstitem
         This primitive is suitable for all instruments.
      }
   }
   \sstdiytopic{
      Tasks
   }{
      KAPPA: STATS.
   }
}
\newpage
\sstroutine{
   CONVERT\_TO\_FITS
}{
   Converts current observation to FITS
}{
   \sstdescription{
      Converts current observation to FITS
   }
}
\newpage
\sstroutine{
   \_CREATE\_GRAPHIC\_FROM\_FILE\_
}{
   Create a PNG, GIF, or JPG graphic from a given file
}{
   \sstdescription{
      This primitive creates a PNG, GIF, or JPG graphic from the
      supplied file. It currently only supports 1-D and 2-D files.
   }
   \sstnotes{
      \sstitemlist{

         \sstitem
         This primitive is suitable only for 1-D and 2-D input files.

         \sstitem
         The graphic file will have the same root filename as the input
         file, with the appropriate file extension.
      }
   }
}
\newpage
\sstroutine{
   \_CREATE\_IMAGE\_CATALOGUE\_\_
}{
   source extraction and photometry on all sources
}{
   \sstdescription{
      For the group file, find all the sources and calculate the flux of
      each detected source. Write the results to a catalogue file in the 
      SExtractor FITS1.0 format.
   }
   \sstnotes{
      \sstitemlist{

         \sstitem
         Currently uses SEXTRACTOR for source extraction and for
         photometry.
      }
   }
   \sstdiytopic{
      Required Perl Modules
   }{
      Starlink::Extractor, Astro::WaveBand, Astro::Catalog.
   }
}
\newpage
\sstroutine{
   \_CREATE\_RAW\_FRAME\_
}{
   Creates a raw frame in ORAC\_DATA\_OUT
}{
   \sstdescription{
      Null primitive for LCOSBIG.
   }
   \sstnotes{
      \sstitemlist{

         \sstitem
         This primitive is suitable for LCOSBIG.
      }
   }
}
\newpage
\newpage
\sstroutine{
   \_DATA\_QC\_TEST\_
}{
   Evaluates the image data quality control flag and sets the
      relevant bitmask keywords in the data catalogue product
}{
   \sstdescription{
      Evaluates the image data quality control flag and sets the
      relevant bitmask keywords in the data catalogue product.
   }
   \sstnotes{
      \sstitemlist{

         \sstitem
            This primitive is suitable for imaging cameras.

         \sstitem
            Processing only occurs for object frames.
      }
   }
   \sstdiytopic{
      Output Data
   }{
      The results of the comparison are returned in
      \$\_DATA\_QC\_TEST\_\{QCPARAM\} and will contain either {\tt '}T{\tt '} if the
      comparison passed, {\tt '}F{\tt '} if the comparison failed or {\tt '}U{\tt '} if the
      comparison couldn{\tt '}t be made.
   }
   \sstdiytopic{
      Output Files
   }{
      None.
   }
}
\newpage
\sstroutine{
   \_DELETE\_RAW\_FRAME\_
}{
   Remove the raw frame file
}{
   \sstdescription{
      Generic primitive to remove the actual raw frame file.
   }
   \sstnotes{
      In rare cases you do not care about the Frame data product (maybe
      because the group product is the only product of interest) and do
      not want it to remain on disk. This primitive will erase the
      current frame file.
   }
}
\newpage
\sstroutine{
   \_DERIVED\_PRODUCTS\_STEER\_
}{
   Steers processing for DERIVED\_PRODUCTS recipe
}{
   \sstdescription{
      This primitive control processing for DERIVED\_PRODUCTS recipe
      through steering headers listed below.
   }
   \sstnotes{
      \sstitemlist{

         \sstitem
         This primitive is suitable for LCOGT imaging CCD cameras.

         \sstitem
         Processing only occurs for object frames.
      }
   }
   \sstdiytopic{
      Steering Headers
   }{
      ADD\_ASTROMETRY = LOGICAL
         Whether or not automated astrometry is to occur. This equates
         to argument DOASTROM.
      CREATE\_CATALOGUE = LOGICAL
         Whether or not catalogue creation is to occur. This equates to
         argument DOIMGCAT.
      CALCULATE\_SEEING\_STATS = LOGICAL
         Whether or not seeing statistics should be calculated. This
         equates to argument DOSEEING.
      QUICK\_LOOK = LOGICAL
         Whether or not to perform quicklook processing. This equates to
         argument DOQUICKLOOK.
   }
}
\newpage
\sstroutine{
   \_DERIVED\_PRODUCTS\_
}{
   Create derived products from the processed frames
}{
   \sstdescription{
      This primitive creates the derived products from the processed
      frames. The following items are performed:

      \sstitemlist{

         \sstitem
         headers are updated with the details of the master calibration
         frames and correction steps used,

         \sstitem
         the saturation and trimmed section are updated with the values
         used in the pipeline processing,

         \sstitem
         A clipped mean, median and sigma are calculated and added to
         the header,

         \sstitem
         the QC modules is run to produce the four QC flags in the
         headers,

         \sstitem
         incorrect values of the TAGID and/or PROPID are flagged and
         corrected and the PROPID is checked against the lists of known
         proposals,

         \sstitem
         access rights and state of the data are set,

         \sstitem
         a 512x512 PNG thumbnail of the processed image is created,

         \sstitem
         the frame is converted to FITS format

         \sstitem
         the pytimecorrect.py code is called to compute per-star
         barycentric time correction and airmass and updates the header of
         the BCD image and source catalog,

         \sstitem
         if a final destination is defined by the enviroment variable
         FINAL\_DATA\_OUT, the derived products are moved there.
      }
   }
   \sstnotes{
      \sstitemlist{

         \sstitem
            This primitive is suitable for imaging cameras.

         \sstitem
            Processing only occurs for object frames.

         \sstitem
            Various incorrect/invalid values of the PROPID and TAGID are
            corrected and overwritten.

         \sstitem
            Data is normally set to private with a publication date of 1
            year from the frames{\tt '} date. Data with a PROPID beginning with
            {\tt "}ENG{\tt "} or Engineering is set to private indefinitely. EPO with a
            PROPID or TAGID beginning with LCOEPO$|$FTPEPO$|$HAWEPO is set to
            public immediate.
      }
   }
   \sstdiytopic{
      External Tasks
   }{
      The following external tasks are used:
      KAPPA
         FITSMOD
      ORAC-DR PRIMITIVES
         \_ADD\_PIPELINE\_VERSION\_, \_CHECK\_PROPID\_, \_CLIPPED\_STATS\_MEDIAN\_,
         \_CONVERT\_TO\_FITS\_, \_CREATE\_GRAPHIC\_FROM\_FILE\_,
         \_GET\_SATURATION\_LEVEL\_, \_SET\_FILE\_FITS\_ITEM\_, \_SET\_QC\_FLAGS\_
   }
   \sstdiytopic{
      Output Files
   }{
      Depending on whether the Frame uhdr entry QUICK\_LOOK is set to 1
      or not, the 00\_bp\_ff suffix of the processed SDF frame is replaced
      with a 10.fits or 90.fits when converting to the output FITS
      frame. The processed SDF frame is also converted to a PNG image
      with a 10.png or 90.png suffix, depending on whether the Frame
      uhdr entry QUICK\_LOOK.
   }
}
\newpage
\sstroutine{
   \_DIVIDE\_BY\_FLAT\_
}{
   Flat-fields a frame
}{
   \sstdescription{
      This primitive divides the current frame by the most-recent and
      matching flat-field frame from \$Cal-$>$flat method.
   }
   \sstnotes{
      \sstitemlist{

         \sstitem
         This primitive is suitable for UFTI, IRCAM, and Michelle in
         imaging mode.

         \sstitem
         Processing only occurs for object and sky frames.

         \sstitem
         The flat-fielded image can be displayed.

         \sstitem
         The frame title is propagated.
      }
   }
   \sstdiytopic{
      Output Data
   }{
      \sstitemlist{

         \sstitem
         Flat-fielded frame inheriting the frame{\tt '}s name but with the \_ff
         suffix.
      }
   }
   \sstdiytopic{
      Tasks
   }{
      CCDPACK: FLATCOR.
   }
}
\newpage
\sstroutine{
   \_FILE\_LIBRARY\_BIAS\_
}{
   Makes a copy of the masterbias and index file in the calibration
   library
}{
   \sstdescription{
      This primitive makes a copy of the masterbias and index.bias from
      the current directory to calibration library (pointed to by the
      ORAC\_DATA\_CAL environment variable.
   }
   \sstnotes{
      \sstitemlist{

         \sstitem
         This primitive is suitable for optical imaging instruments.

         \sstitem
         Processing only occurs for bias frames, and when the steering
         header MAKE\_BIAS is true.
      }
   }
   \sstdiytopic{
      Output Data
   }{
      The masterbias. It is called \texttt{"} bias\_$<$instrument$>$\_$<$UT
      night$>$\_bin$<$Xbin$>$x$<$Ybin$>$\texttt{"}  where $<$instrument$>$ is the LCOGT
      instrument identifier and $<$Xbin$>$, $<$Ybin$>$ are the binning factors
      in X and Y.
   }
   \sstdiytopic{
      Tasks
   }{
      KAPPA: NDFCOPY.
   }
}
\newpage
\sstroutine{
   \_FILE\_LIBRARY\_DARK\_
}{
   Makes a copy of the masterdark and index file in the calibration
   library
}{
   \sstdescription{
      This primitive makes a copy of the masterdark and index.dark from
      the current directory to calibration library (pointed to by the
      ORAC\_DATA\_CAL environment variable.
   }
   \sstnotes{
      \sstitemlist{

         \sstitem
         This primitive is suitable for optical imaging instruments.

         \sstitem
         Processing only occurs for dark frames, and when the steering
         header MAKE\_DARK is true.
      }
   }
   \sstdiytopic{
      Output Data
   }{
      The masterdark. It is called \texttt{"} dark\_$<$instrument$>$\_$<$UT
      night$>$\_bin$<$Xbin$>$x$<$Ybin$>$\texttt{"}  where $<$instrument$>$ is the LCOGT
      instrument identifier and $<$Xbin$>$, $<$Ybin$>$ are the binning factors
      in X and Y.
   }
   \sstdiytopic{
      Tasks
   }{
      KAPPA: NDFCOPY.
   }
}
\newpage
\sstroutine{
   \_FILE\_LIBRARY\_FLAT\_
}{
   Makes a copy of the masterflat and index file in the calibration
   library
}{
   \sstdescription{
      This primitive makes a copy of the masterflat and index.flat from
      the current directory to calibration library (pointed to by the
      ORAC\_DATA\_CAL environment variable.
   }
   \sstnotes{
      \sstitemlist{

         \sstitem
         This primitive is suitable for optical imaging instruments.

         \sstitem
         Processing only occurs for flat frames, and when the steering
         header MAKE\_FLAT is true.
      }
   }
   \sstdiytopic{
      Output Data
   }{
      The masterflat. It is called
      \texttt{"} flat\_$<$instrument$>$\_$<$UTnight$>$\_SKYFLAT\_bin$<$Xbin$>$x$<$Ybin$>$\_$<$filter$>$\texttt{"}
      where $<$instrument$>$ is the LCOGT instrument identifier, $<$Xbin$>$,
      $<$Ybin$>$ are the binning factors in X and Y and $<$filter$>$ is the
      filter.
   }
   \sstdiytopic{
      Tasks
   }{
      KAPPA: NDFCOPY.
   }
}
\newpage
\sstroutine{
   \_FLAT\_QC\_
}{
   To perform Quality Control on flat fields produced by this
   pipeline
}{
   \sstdescription{
      This primitive looks for the separate master flat fields produces
      in each filter at both the start and the end of a single night of
      data. It performs the following tests:

      \sstitemlist{

         \sstitem
         If morning and evening twilight master flats were produced:

         \sstitem
            Compare the am/pm masters and check the RMS of the residuals is
            within bounds (implementation of J. Eastman{\tt '}s IDL algorithm).
      }
   }
   \sstnotes{
      This is an ORAC-DR/Perl implementation of J. Eastman{\tt '}s IDL
      algorithm.
   }
   \sstdiytopic{
      External Tasks
   }{
      KAPPA:
         STATS, SUB
      ORAC-DR PRIMITIVES
         \_DELETE\_A\_FRAME\_
   }
   \sstdiytopic{
      Output Data
   }{
      None.
   }
}
\newpage
\sstroutine{
   \_GET\_FILTER\_PARAMETERS\_
}{
   Returns LCOGT optical filter characteristics
}{
   \sstdescription{
      This primitive using a switch structure to return two
      characteristics or relating to the LCOGT optical imaging filters
      through arguments. Default values are returned if the filter is
      not recognised.
   }
   \sstnotes{
      \sstitemlist{

         \sstitem
         This primitive is suitable for LCOGT optical photometry.

         \sstitem
         The filter name comes from the user header ORAC\_FILTER.

         \sstitem
         The recognised filters have names ending with U, B, V, R, I,
         up, gp, rp, ip, zp, zs and Y.

         \sstitem
         The mean extinction coefficients are: U: 0.53, B: 0.27, V:
         0.14, R: 0.10, I: 0.05, u{\tt '}: 0.56, g{\tt '}: 0.20, r{\tt '}:0.11, i{\tt '}: 0.05, z{\tt '}:
         0.04 and Y: 0.03. Zero extinction applies to any other filter.
         Effective wavelengths are from Bessell (1998); extinction is from
         AJP{\tt '}s loc\_m1\_assm.c code as of 2012/05/08 for ELP; M1 inst.
         zeropoints are from AJP{\tt '}s loc\_m1\_assm.c code as of 2012/05/08 for
         1m0 kb72.
      }
   }
}
\newpage
\sstroutine{
   \_GET\_GAIN\_
}{
   Finds the LCO instrument gain in electrons per ADU for the current
   Frame
}{
   \sstdescription{
      This primitive obtains the gain in electrons per ADU for the
      current frame. It first attempts to find a value from the
      ORAC\_GAIN header. If this is null or less than 3.0, the primitive
      uses a time-dependent default value, and it reports the use of the
      default.

      The gain is returned through an argument.
   }
   \sstnotes{
      \sstitemlist{

         \sstitem
         This primitive is suitable for LCO optical instruments.
      }
   }
}
\newpage
\sstroutine{
   \_GET\_READNOISE\_
}{
   Finds the LCO instrument readnoise in electrons for the current
   Frame
}{
   \sstdescription{
      This primitive obtains the readnoise in electrons for the current
      frame. It first attempts to find a value for the calibration
      system. Where there is no value, it tries a header for the value,
      and if that{\tt '}s not defined, the primitive assigns a default.

      The readnoise is returned through an argument.
   }
   \sstnotes{
      \sstitemlist{

         \sstitem
         This primitive is suitable for LCO optical instruments.

         \sstitem
         The read noise comes from the readnoise calibration, or failing
         that the header RDNOISE.
      }
   }
}
\newpage
\sstroutine{
   \_GET\_SATURATION\_LEVEL\_
}{
   Finds the LCOSBIG saturation level in ADU the current Frame
}{
   \sstdescription{
      This primitive obtains the LCOSBIG saturation level in ADU for the
      current frame.

      The saturation level is returned through an argument.
   }
   \sstnotes{
      \sstitemlist{

         \sstitem
         This primitive is only suitable for LCOSBIG.

         \sstitem
         The values are estimated from fullwell and gain on the
         instrument{\tt '}s Web page, viz.
         http://www.sbig.com/large\_format/6303E\_specs.htm.
      }
   }
}
\newpage
\newpage
\newpage
\sstroutine{
   \_INSTRUMENT\_HELLO\_
}{
   Performs the LCOSBIG-specific imaging setup
}{
   \sstdescription{
      This primitive is performs the instrument specific setup for
      imaging. It{\tt '}s needed for the generic \_IMAGING\_HELLO\_. In this case
      it merely reports that the set-up operations are complete.
   }
   \sstnotes{
      \sstitemlist{

         \sstitem
         This primitive is suitable for LCOSBIG.
      }
   }
}
\newpage
\sstroutine{
   \_LCOGT\_STANDARD\_MAGNITUDE\_
}{
   Obtains the catalogue magnitude of a LCOGT standard
}{
   \sstdescription{
      This primitive reads the faint-standard catalogue or its
      predecessor. A case- and space-insensitive comparison of the
      supplied object name with the entries in the table provides a
      catalogue magnitude in U, B, V, R, or I for a standard star.
   }
   \sstnotes{
      \sstitemlist{

         \sstitem
         This primitive is suitable for LCOGT optical imagers.

         \sstitem
         Processing only occurs when it is time to perform photometry,
         i.e. when the steering header DO\_APHOT is true.

         \sstitem
         An error occurs when the filter is not one of U, B, V, R, or I.

         \sstitem
         The standard-star catalogue used is
         \$ORAC\_DATA\_CAL/landolt\_ubvri.dat. An error results when the
         catalogue cannot be opened.
      }
   }
}
\newpage
\sstroutine{
   \_MAKE\_BIAS\_FROM\_GROUP\_
}{
   Makes a masterbias from the current group of frames
}{
   \sstdescription{
      This primitive makes a master bias from the current group. The
      primitive files the resultant bias in its calibration index.
   }
   \sstnotes{
      \sstitemlist{

         \sstitem
         This primitive is suitable for optical imaging instruments.

         \sstitem
         Processing only occurs for bias frames, and when the steering
         header MAKE\_BIAS is true.

         \sstitem
         The bias is displayed.
      }
   }
   \sstdiytopic{
      Output Data
   }{
      The masterbias. It is called {\tt "}bias\_$<$instrument$>$\_$<$UT
      night$>$\_bin$<$Xbin$>$x$<$Ybin$>${\tt "} where $<$instrument$>$ is the LCOGT
      instrument identifier and $<$Xbin$>$, $<$Ybin$>$ are the binning factors
      in X and Y.
   }
   \sstdiytopic{
      Tasks
   }{
      CCDPACK: MAKEBIAS; KAPPA: FITSMOD, NDFCOPY.
   }
}
\newpage
\sstroutine{
   \_MAKE\_DARK\_FROM\_GROUP\_
}{
   Makes a masterdark from the current group of frames
}{
   \sstdescription{
      This primitive makes a master dark from the current group. The
      primitive files the resultant dark in its calibration index.
   }
   \sstnotes{
      \sstitemlist{

         \sstitem
         This primitive is suitable for optical imaging instruments.

         \sstitem
         Processing only occurs for dark frames, and when the steering
         header MAKE\_DARK is true.

         \sstitem
         The dark is displayed.
      }
   }
   \sstdiytopic{
      Output Data
   }{
      The masterdark. It is called {\tt "}dark\_$<$instrument$>$\_$<$UT
      night$>$\_bin$<$Xbin$>$x$<$Ybin$>${\tt "} where $<$instrument$>$ is the LCOGT
      instrument identifier and $<$Xbin$>$, $<$Ybin$>$ are the binning factors
      in X and Y.
   }
   \sstdiytopic{
      Tasks
   }{
      CCDPACK: MAKEDARK; KAPPA: FITSMOD, NDFCOPY.
   }
}
\newpage
\sstroutine{
   \_MAKE\_FLAT\_FROM\_GROUP\_
}{
   Makes flats from the current group of frames
}{
   \sstdescription{
      This primitive makes self flats from the current group, one for
      each distinct observation filter. For eacg flat it uses a median
      to combine the frames pixel by pixel, and then divides the
      resultant image by its mean form the flat field. The primitive
      files the resultant flat in its calibration index.
   }
   \sstnotes{
      \sstitemlist{

         \sstitem
         This primitive is suitable for infrared imaging instruments.

         \sstitem
         Processing only occurs for object, sky, or calibration-lamp
         frames, and when the steering header MAKE\_FLAT is true.

         \sstitem
         The list of filters present in the group is listed in an array
         stored by reference in the group user header FILTER\_LIST. If this
         is undefined, only a single flat is made for filter stored in the
         current Frame{\tt '}s user header ORAC\_FILTER.

         \sstitem
         There is special behaviour for a combined polarimetry flat (see
         {\tt "}OUTPUT DATA{\tt "}). The string {\tt "}pol{\tt "} in the filter name is used to
         indicate polarimetry data.

         \sstitem
         The flat is displayed.
      }
   }
   \sstdiytopic{
      Output Data
   }{
      The flat field. It is called is {\tt "}flat\_$<$filter$>$\_$<$groupnumber$>${\tt "} for
      the first cycle, and {\tt "}flat\_$<$filter$>$\_$<$groupnumber$>$\_c$<$cyclenumber$>${\tt "}
      for subsequent cycles, where $<$groupnumber$>$ is the frame number of
      the group, $<$filter$>$ is the filter name, and $<$cyclenumber$>$ is the
      cycle number derived from steering header CYCLE\_NUMBER.
      An exception is for polarimetric data, where the name becomes
      flat\_$<$filter$>$\_pol$<$waveplate\_angle$>$\_$<$groupnumber$>$. The
      $<$waveplate\_angle$>$ is the integer part of the angle, e.g. 22, 67,
      from internal header ORAC\_WAVEPLATE\_ANGLE. Subsequent cycles for
      polarimetry also have the {\tt "}\_c$<$cyclenumber$>${\tt "} suffix, but the cycle
      comes from steering header POL\_CYCLE\_NUMBER. When steering header
      WAVEPLATE\_FLAT is false (0), copies of the flat are made, one for
      each angle, using the above nomenclature. Each has its waveplate
      angle set to its nominal angle. This allows a single
      ORAC\_WAVEPLATE\_ANGLE rule entry irrespective of whether all
      waveplate angles were combined to make a flat or not.
   }
   \sstdiytopic{
      Tasks
   }{
      CCDPACK: MAKEFLAT; KAPPA: FITSMOD, NDFCOPY.
   }
}
\newpage
\sstroutine{
   \_NIGHT\_LOG\_
}{
   Produces a text file log of a night{\tt '}s imaging observations
}{
   \sstdescription{
      This recipe takes a night{\tt '}s imaging observations, and creates a
      text file containing a headed tabulation of parameters for each
      frame.

      The parameters are: observation number, group number, object name,
      observation type, UT start time, exposure time, number of coadds,
      read mode and speed, filter, start airmass, frame dimensions in
      pixels, base equatorial co-ordinates, and data-reduction recipe
      name.
   }
   \sstnotes{
      \sstitemlist{

         \sstitem
         The $<$date$>$ comes from the internal header keyword ORAC\_UTDATE.

         \sstitem
         The logfile created by this primitive does not follow the
         standard ORAC-DR naming convention (log.xxxx) since it can be used
         to write log files to directories other than \$ORAC\_DATA\_OUT and
         unique file names are required.

         \sstitem
         Fudges missing or old headers.

         \sstitem
         Uses user header ORAC\_INSTRUMENT to specify the file name.

         \sstitem
         Specification provided by Sandy Leggett.
      }
   }
   \sstdiytopic{
      Output Data
   }{
      \sstitemlist{

         \sstitem
         The text log file \$ORAC\_DATA\_IN/$<$date$<$.nightlog, where $<$date$>$
         is the UT date, unless the OUT argument is set, whereupon the log
         is in \$ORAC\_DATA\_OUT. This enables a separate on-the-fly log. For
         the multi-mode instruments UIST, Michelle, IRIS2, and ISAAC the
         file is \$ORAC\_DATA\_IN/$<$date$>$\_im.nightlog.
         For LCOGT data, we make use of the DAY-OBS header keyword to use
         as the observation date (and hence log root) to prevent writing
         two logs after we go over the UTC midnight boundary.
         The on-the-fly log in \$ORAC\_DATA\_OUT is always appended to, being
         created only if it doesn{\tt '}t exist. Thus multiple entries for the
         same observation may exist in the on-the-fly log if the pipeline
         is rerun.
         The {\tt "}clean{\tt "} log file in F\$$<$ORAC\_DATA\_IN$>$ is re-started if the
         observation number equals 1 and is appended to otherwise, being
         created as necessary.
      }
   }
}
\newpage
\sstroutine{
   \_OFFLINE\_REDUCTION\_HELLO\_
}{
   Sets up data-reduction tasks and data for OFFLINE\_REDUCTION
   recipes
}{
   \sstdescription{
      Sets up CCDPACK-related global parameters for OFFLINE\_REDUCTION
      recipes. The settings are as follows.

      \sstitemlist{

         \sstitem
         The readout bounds in the internal headers ORAC\_X\_LOWER\_BOUND,
         ORAC\_Y\_LOWER\_BOUND, ORAC\_X\_UPPER\_BOUND, ORAC\_Y\_UPPER\_BOUND define
         the pixel limits for processing, i.e. there are no bias strips and
         interpolation direction.

         \sstitem
         Error processing is disabled so the readout noise and
         analogue-to-digital conversions are not specified.

         \sstitem
         There is no deferred charge.

         \sstitem
         Position list processing tasks expect to find the names of
         lists stored within NDFs.

         \sstitem
         Logging is to the terminal.

         \sstitem
         The data type of NDF arrays is preserved.

         \sstitem
         Does not detect saturated pixels.

         \sstitem
         Parameters are neither saved from or to a `restoration{\tt '} file.

      }
      The script also performs the following tasks.

      \sstitemlist{

         \sstitem
         Calls the steering primitive to set steering headers.

         \sstitem
      }
   }
   \sstnotes{
      \sstitemlist{

         \sstitem
         This primitive is suitable for imaging instruments.
      }
   }
   \sstdiytopic{
      Tasks
   }{
      CCDPACK: CCDSETUP.
   }
}
\newpage
\sstroutine{
   \_QUICK\_LOOK\_HELLO\_
}{
   Sets up data-reduction tasks and data for QUICK\_LOOK recipes
}{
   \sstdescription{
      Sets up CCDPACK-related global parameters for QUICK\_LOOK recipes.
      The settings are as follows.

      \sstitemlist{

         \sstitem
         The readout bounds in the internal headers ORAC\_X\_LOWER\_BOUND,
         ORAC\_Y\_LOWER\_BOUND, ORAC\_X\_UPPER\_BOUND, ORAC\_Y\_UPPER\_BOUND define
         the pixel limits for processing, i.e. there are no bias strips and
         interpolation direction.

         \sstitem
         Error processing is disabled so the readout noise and
         analogue-to-digital conversions are not specified.

         \sstitem
         There is no deferred charge.

         \sstitem
         Position list processing tasks expect to find the names of
         lists stored within NDFs.

         \sstitem
         Logging is to the terminal.

         \sstitem
         The data type of NDF arrays is preserved.

         \sstitem
         Does not detect saturated pixels.

         \sstitem
         Parameters are neither saved from or to a `restoration{\tt '} file.

      }
      The script also performs the following tasks.

      \sstitemlist{

         \sstitem
         Calls the steering primitive to set steering headers.

         \sstitem
      }
   }
   \sstnotes{
      \sstitemlist{

         \sstitem
         This primitive is suitable for imaging instruments.
      }
   }
   \sstdiytopic{
      Tasks
   }{
      CCDPACK: CCDSETUP.
   }
}
\newpage
\sstroutine{
   \_REDUCE\_BIAS\_HELLO\_
}{
   Sets up data-reduction tasks and data for REDUCE\_BIAS recipes
}{
   \sstdescription{
      This primitive sets up CCDPACK-related global parameters for the
      REDUCE\_BIAS recipe and performs preliminary data reduction.

      The CCDPACK settings are as follows.

      \sstitemlist{

         \sstitem
         The readout bounds in the internal headers ORAC\_X\_LOWER\_BOUND,
         ORAC\_Y\_LOWER\_BOUND, ORAC\_X\_UPPER\_BOUND, ORAC\_Y\_UPPER\_BOUND define
         the pixel limits for processing, i.e. there are no bias strips and
         interpolation direction.

         \sstitem
         Error processing is disabled so the readout noise and
         analogue-to-digital conversions are not specified.

         \sstitem
         There is no deferred charge.

         \sstitem
         Position list processing tasks expect to find the names of
         lists stored within NDFs.

         \sstitem
         Logging is to the terminal.

         \sstitem
         The data type of NDF arrays is preserved.

         \sstitem
         Does not detect saturated pixels.

         \sstitem
         Parameters are neither saved from or to a `restoration{\tt '} file.

      }
      The script also performs the following tasks.

      \sstitemlist{

         \sstitem
         Calls the steering primitive to set steering headers.

         \sstitem
         Calls \_MASK\_BAD\_PIXELS\_ primitive to mask bad pixels.
      }
   }
   \sstnotes{
      \sstitemlist{

         \sstitem
         This primitive is suitable for imaging instruments.
      }
   }
   \sstdiytopic{
      Steering Headers
   }{
      USE\_VARIANCE = LOGICAL
         Whether or not variance processing is to occur. This equates to
         argument USEVAR.
   }
}
\newpage
\sstroutine{
   \_REDUCE\_BIAS\_STEER\_
}{
   Steers processing for REDUCE\_BIAS recipe
}{
   \sstdescription{
      This primitive control processing for REDUCE\_BIAS recipe through
      steering headers listed below.
   }
   \sstnotes{
      \sstitemlist{

         \sstitem
         This primitive is suitable for imaging optical cameras.

         \sstitem
         Processing only occurs for bias frames.
      }
   }
   \sstdiytopic{
      Steering Headers
   }{
      BIAS\_FRAMES = HASH
         The name of the bias for a given exposure time, and the
         corresponding number of frames used to create it.
      USE\_VARIANCE = LOGICAL
         Whether or not variance processing is to occur. This equates to
         argument USEVAR.
   }
}
\newpage
\sstroutine{
   \_REDUCE\_BIAS\_TIDY\_
}{
   Removes unwanted intermediate files for the REDUCE\_BIAS recipe
}{
   \sstdescription{
      Removes all intermediate frames. Files are only removed when they
      are no longer needed, as guided by the steering header MAKE\_BIAS.
   }
}
\newpage
\sstroutine{
   \_REDUCE\_DARK\_STEER\_
}{
   Steers processing for REDUCE\_DARK recipe
}{
   \sstdescription{
      This primitive control processing for REDUCE\_DARK recipe through
      steering headers listed below.
   }
   \sstnotes{
      \sstitemlist{

         \sstitem
         This primitive is suitable for imaging optical cameras.

         \sstitem
         Processing only occurs for dark frames.
      }
   }
   \sstdiytopic{
      Steering Headers
   }{
      USE\_VARIANCE = LOGICAL
         Whether or not variance processing is to occur. This equates to
         argument USEVAR.
   }
}
\newpage
\sstroutine{
   \_REDUCE\_DARK\_TIDY\_
}{
   Removes unwanted intermediate files for the REDUCE\_DARK recipe
}{
   \sstdescription{
      Removes all intermediate frames. Files are only removed when they
      are no longer needed, as guided by the steering header MAKE\_DARK.
   }
}
\newpage
\sstroutine{
   \_REMOVE\_BIAS\_
}{
   Subtracts a bias frame
}{
   \sstdescription{
      This primitive subtracts a zero bias from the current frame;
      unless the data have variance information and were taken using a
      non-ND mode (i.e. where the bias has not already been subtracted),
      whereupon a bias frame, if available, is subtracted.

      For most instruments there is no bias to subtract so it is
      something of a placeholder primitive. Its main purpose is to set
      up CCDPACK for subsequent processing. For instance, CCDPACK will
      complain if debiassing is not performed before say flat-fielding.
      The primitive reports a successful bias subtraction and the frames
      concerned.
   }
   \sstnotes{
      \sstitemlist{

         \sstitem
         This primitive is suitable for UFTI, IRCAM, INGRID, and
         Michelle in imaging mode.

         \sstitem
         Processing occurs for all frames, and sub-frames therein.

         \sstitem
         Where a bias frame is used, it is the most-recent and matching
         given by \$Cal-$>$bias method.

         \sstitem
         The observing mode (read type) comes from user header
         ORAC\_DETECTOR\_READ\_TYPE.

         \sstitem
         The bias-subtracted image can be displayed.

         \sstitem
         The frame title is propagated.
      }
   }
   \sstdiytopic{
      Output Data
   }{
      \sstitemlist{

         \sstitem
         Bias-subtracted frame inheriting the frame{\tt '}s name but with the
         \_db suffix.
      }
   }
   \sstdiytopic{
      Tasks
   }{
      CCDPACK: DEBIAS.
   }
}
\newpage
\sstroutine{
   \_SET\_ORIGIN\_
}{
   Sets the origin of an observation
}{
   \sstdescription{
      This primitive sets the origin of an observation, including all
      integrations, using the ORAC\_X\_LOWER\_BOUND and ORAC\_Y\_LOWER\_BOUND
      user headers in the frame. It is needed to correct the raw data
      from a sub-array for which the origin is still at the default. If
      either header is undefined, the primitive creates a default origin
      (1,1). The origin is not set if ORAC\_X\_LOWER\_BOUND does not exist.
   }
   \sstnotes{
      \sstitemlist{

         \sstitem
         This primitive is suitable for UFTI, IRCAM, and Michelle and
         UIST in imaging mode.
      }
   }
   \sstdiytopic{
      Tasks
   }{
      KAPPA: SETORIGIN.
   }
}
\newpage
\newpage
\sstroutine{
   \_SKY\_FLAT\_STEER\_
}{
   Steers processing for SKY\_FLAT recipes
}{
   \sstdescription{
      This primitive control processing for SKY\_FLAT recipes through
      steering headers listed below.
   }
   \sstnotes{
      \sstitemlist{

         \sstitem
         This primitive is suitable for imaging infrared cameras.

         \sstitem
         Processing only occurs for object, sky, or calibration lamp
         frames.

         \sstitem
         The data are deemed to be polarimetry if the frame internal
         header ORAC\_FILTER contains the string {\tt "}pol{\tt "}.

         \sstitem
         A list of the distinct filters within the group is stored in an
         array stored by reference in the group user header FILTER\_LIST.
      }
   }
   \sstdiytopic{
      Steering Headers
   }{
      CYCLE\_NUMBER = INTEGER
         Number of the cycle, a cycle being a set of frames to complete
         a pass through the recipe. The first cycle is 0.
      JITTER\_NUMBER = INTEGER
         The number of frames in the jitter.
      MAKE\_FLAT = LOGICAL
         Whether or not to make the flat. The flat is made once all the
         jittered target frames in a cycle are available.
      MASK\_OBJECTS = LOGICAL
         Whether or not to mask the objects. Masking occurs when all the
         jittered frames in a cycle are available.
      POL\_CYCLE\_NUMBER = INTEGER
         Number of the polarimetry cycle, a cycle being a set of frames
         to complete a pass through the recipe for all waveplate angles.
         The first cycle is 0.
      USE\_VARIANCE = LOGICAL
         Whether or not variance processing is to occur. This equates to
         argument USEVAR.
      WAVEPLATE\_FLAT = LOGICAL
         See the argument of the same name. This header merely
         propagates the value of the argument.
   }
}
\newpage
\sstroutine{
   \_STANDARD\_MAGNITUDE\_
}{
   Obtains the catalogue magnitude of a standard
}{
   \sstdescription{
      This primitive reads the faint-standard catalogue or its
      predecessor. A case- and space-insensitive comparison of the
      supplied object name with the entries in the table provides a
      catalogue magnitude in U, B, V, R, or I for a standard star.
   }
   \sstnotes{
      \sstitemlist{

         \sstitem
         This primitive is suitable for LCOGT optical imagers.

         \sstitem
         Processing only occurs when it is time to perform photometry,
         i.e. when the steering header DO\_APHOT is true.

         \sstitem
         An error occurs when the filter is not one of U, B, V, R, or I.

         \sstitem
         Invokes \_LCOGT\_STANDARD\_MAGNITUDE\_ to obtain the magnitude and
         catalogue name.
      }
   }
}
\newpage
\sstroutine{
   \_SUBTRACT\_DARK\_NO\_THRESH\_
}{
   Subtracts a dark frame
}{
   \sstdescription{
      This primitive subtracts from the current frame the most-recent
      and matching dark frame given by \$Cal-$>$dark method. It reports a
      successful dark subtraction and the frames concerned.
   }
   \sstnotes{
      \sstitemlist{

         \sstitem
         This primitive is suitable for UFTI, IRCAM, and Michelle in
         imaging mode.

         \sstitem
         Processing only occurs for object and sky frames.

         \sstitem
         The dark-subtracted image can be displayed.

         \sstitem
         The subtraction assumes the same exposure time for the dark and
         object frame. That validation should be done by the \$Cal-$>$dark
         method.

         \sstitem
         The frame title is propagated.
      }
   }
   \sstdiytopic{
      Output Data
   }{
      \sstitemlist{

         \sstitem
         Dark-subtracted frame inheriting the frame{\tt '}s name but with the
         \_dk suffix.
      }
   }
   \sstdiytopic{
      Tasks
   }{
      CCDPACK: CALCOR.
   }
}
\newpage
\sstroutine{
   \_SUBTRACT\_DARK\_
}{
   Subtracts a dark frame
}{
   \sstdescription{
      This primitive subtracts from the current frame the most-recent
      and matching dark frame given by \$Cal-$>$dark method. It reports a
      successful dark subtraction and the frames concerned.

      Since transient `hot{\tt '} and `cold{\tt '} pixels can be present despite the
      application of a bad-pixel mask, the primitive also thresholds the
      dark-subtracted frame, setting values beyond the limits to be bad
      (i.e. undefined), to remove these non-physical values. Such values
      can lead to problems later in the pipeline. In a sense this
      processing step augments the bad-pixel mask.
   }
   \sstnotes{
      \sstitemlist{

         \sstitem
         This primitive is suitable for infrared imaging.

         \sstitem
         Processing only occurs for object, sky, and flat frames.

         \sstitem
         The dark-subtracted image can be displayed.

         \sstitem
         The subtraction assumes the same exposure time for the dark and
         object frame. That validation should be done by the \$Cal-$>$dark
         method.

         \sstitem
         The lower threshold limit is the clipped mean (mode) minus five
         standard deviations, subject to the constraint that the limit lies
         between -100 and 1. The upper limit is 1000 above the nominal
         saturation level for the instrument and its mode.

         \sstitem
         The primitive issues a warning if the dark-subtracted frame{\tt '}s
         mode is negative, allowing for the error of the mode. It aborts
         with an error message if the modal dark-subtracted signal is more
         than one standard deviation negative.

         \sstitem
         The frame title is propagated.
      }
   }
   \sstdiytopic{
      Output Data
   }{
      \sstitemlist{

         \sstitem
         Dark-subtracted frame inheriting the frame{\tt '}s name but with the
         \_dk suffix.

         \sstitem
         An array with bad-value substitution beyond thresholds and
         inheriting the frame{\tt '}s name but with the \_th suffix.
      }
   }
   \sstdiytopic{
      Tasks
   }{
      CCDPACK: CALCOR; KAPPA: THRESH.
   }
}
\newpage
\sstroutine{
   \_TURN\_ON\_HISTORY\_
}{
   Switches on history recording
}{
   \sstdescription{
      This primitive enables NDF history recording for each integration
      in an observation.
   }
   \sstnotes{
      \sstitemlist{

         \sstitem
         This primitive is suitable for UFTI, IRCAM, and Michelle and
         UIST in imaging mode.

         \sstitem
         If the ORAC\_HISTORY\_OFF environment variable is set, then
         history will be disabled.
      }
   }
   \sstdiytopic{
      Tasks
   }{
      KAPPA: HISSET.
   }
}


\newpage

\section{\xlabel{ap_qc_indexfiles}Specification of LCOGT Quality Control index files\label{ap:qc_indexfiles}}

The following pages describe the LCOGT Quality Control index files in detail.
These files are part of the calibration scheme therefore they are controlled by
LCOGT-specific subclasses of \texttt{ORAC::Calib} (found in \texttt{<ORAC-DR
root>/src/lib/perl5/ORAC/Calib/)}) namely \texttt{ORAC::Calib::LCOSBIG} (used
for the \textsc{LCOSBIG}, \textsc{LCOSBIG\_0M8} (Sedgwick),
\textsc{LCOSINISTRO}, \textsc{LCOFLI}, \textsc{LCOMEROPE} and 
\textsc{LCOSPECTRAL} instruments),  \texttt{ORAC::Calib::LCOSBIG\_0M4} (used for
\textsc{LCOSBIG\_0M4} (0.4m) instruments) and \texttt{ORAC::Calib::LCOCC} (used
for Context Camera instruments). In addition there is the basis of
low-resolution spectrographic calibration support through
\texttt{ORAC::Calib::LCOFLOYDS}.

The columns to appear in the index file are controlled by a list of column names
within a subroutine of the same name (e.g. the format of the \texttt{index.dqc}
file is controlled by a \texttt{'dqc'} subroutine) and extra columns can be
added as desired. Once defined, the ORAC-DR calibration infrastructure handles
the location, creation, indexing, updating, retrieval, consistency checking and
verification of the index files. In addition to the declared columns, columns
for the frame filename and the time are automatically created.

\subsection{\xlabel{ap_index.astromqc}Specification of the \texttt{index.astromqc} file\label{ap:index.astromqc}}

This index is used to store the result of the astrometric calibration i.e. the
WCS fit. The format of the file is a header of the column names (with a hash
character in the first column) on the first line followed by a line entry for
each frame consisting of multiple colums with the form of \texttt{<FILENAME>
OFFSETDEC OFFSETRA <ORACTIME> SECPIX XRMS YRMS} where:
\begin{itemize}
\item \texttt{<FILENAME>} is the name of the reduced frame, 
\item \texttt{NSTARS} is the number of stars used in the astrometric fit (set to 0 if
no satisfactory fit was determined), 
\item \texttt{OFFSETDEC} is the offset in Dec in arcseconds between the fitted
frame center and the Dec of the telescope (set to 0 if no satisfactory fit was determined), 
\item \texttt{OFFSETRA} is the offset in RA in arcseconds between the fitted
frame center and the RA of the telescope (set to 0 if no satisfactory fit was determined),
\item \texttt{<ORACTIME>} is the auto-inserted time of frame in standard ORAC
YYYYMMDD.ddddd format,
\item \texttt{SECPIX} is the fitted pixel scale in arcseconds/pixel (set to the
nominal pixel scale if no satisfactory fit was determined)
\item \texttt{XRMS} is the RMS of the astrometric fit in the X direction in
arcseconds (set to '-99' if no satisfactory fit was determined)
\item \texttt{YRMS} is the RMS of the astrometric fit in the X direction in
arcseconds (set to '-99' if no satisfactory fit was determined)
\end{itemize}

A section of an example file is shown below:
\begin{terminalv}
#NSTARS OFFSETDEC OFFSETRA ORACTIME SECPIX XRMS YRMS
lsc1m005-kb78-20150124-0065-s00_bp_ff 0 0 0 20150125.0169907 0.467 -99 -99
lsc1m005-kb78-20150124-0066-s00_bp_ff 0 0 0 20150125.0178704 0.467 -99 -99
lsc1m005-kb78-20150124-0067-s00_bp_ff 1 94.7442281467999 -65.0030048559793 20150125.0189352 0.4699 0.00000 0.00000
lsc1m005-kb78-20150124-0068-s00_bp_ff 1 93.4636426392003 -64.9209369759612 20150125.0209722 0.4700 0.00000 0.00000
lsc1m005-kb78-20150124-0069-s00_bp_ff 12 92.9154980000001 -64.8774346599973 20150125.0231713 0.4700 0.11851 0.07415
lsc1m005-kb78-20150124-0070-s00_bp_ff 12 92.8168515195999 -64.9570614880085 20150125.0243056 0.4700 0.11386 0
\end{terminalv}

The index file is updated by filing new astrometric fit results with the
calibration system in the \task{\_ADD\_AUTO\_ASTROMETRY\_} primitive.

\subsection{\xlabel{ap_index.dqc}Specification of the \texttt{index.dqc} file\label{ap:index.dqc}}

This index is used to store the result of the data quality control. The format of the file is a header of the column names (with a hash
character in the first column) on the first line followed by a line entry for
each frame consisting of multiple colums with the form of \texttt{<FILENAME>
AIRMASS ELLIPTICITY FWHM <ORACTIME> ORIENT QC\_OBCON QC\_IMGST QC\_CATST QC\_PHTST} where:
\begin{itemize}
\item \texttt{<FILENAME>} is the name of the reduced frame, 
\item \texttt{AIRMASS} is the mean airmass at the middle of the exposure, 
\item \texttt{ELLIPTICITY} is the mean ellipticity (0..1) of the sources in the
frame determined by \SEx\  (set to -99 if no seeing statistics were determined), 
\item \texttt{FWHM} is the mean Full Width Half Maximum in arcseconds of the
sources in the frame determined by \SEx\  (set to -99 if no seeing statistics were determined),
\item \texttt{<ORACTIME>} is the auto-inserted time of frame in standard ORAC
YYYYMMDD.ddddd format,
\item \texttt{ORIENT} is the mean orientation in degrees of the sources in the
frame fitted by \SEx\  (set to -99' if no seeing statistics were determined)
\item \texttt{QC\_CATST} is the bitmask of flags set by the QC module for the
catalog extraction (see the \SAICD\ or the \SA\ data release notes for the full
list of flags),
\item \texttt{QC\_IMGST} is the bitmask of flags set by the QC module for the
processed image (see the \SAICD\ or the \SA\ data release notes for the full
list of flags),
\item \texttt{QC\_OBCON} is the string concatanation of tests set by the QC
module for the observing constraints (see the \SAICD\ or the \SA\ data release
notes for the full list of flags),
\item \texttt{QC\_PHTST} is the bitmask of flags set by the QC module for the
photometric calibration (see the \SAICD\ or the \SA\ data release notes for the full
list of flags),
\end{itemize}

A section of an example file is shown below:
\begin{terminalv}
#AIRMASS ELLIPTICITY FWHM ORACTIME ORIENT QC_CATST QC_IMGST QC_OBCON QC_PHTST
lsc1m005-kb78-20150124-0061-s00_bp_ff 1.2029843 0.19747 2.746 20150125.0125926 23.140 1 0 UUUUU 4
lsc1m005-kb78-20150124-0062-s00_bp_ff 1.2049946 0.22837 2.723 20150125.0137037 33.649 1 0 UUUUU 4
lsc1m005-kb78-20150124-0063-s00_bp_ff 1.2072453 0.21732 2.737 20150125.0150231 25.667 1 0 UUUUU 4
lsc1m005-kb78-20150124-0064-s00_bp_ff 1.2089885 0.20205 2.628 20150125.0159259 19.913 1 0 UUUUU 4
\end{terminalv}

The index file is updated by filing new seeing statistic results with the
calibration system in the \task{\_CALCULATE\_SEEING\_STATS\_} primitive and the
index entry is updated to set the QC flags by the .\task{\_SET\_QC\_FLAGS\_} primitive

\subsection{\xlabel{ap_index.zeropoint}Specification of the \texttt{index.zeropoint} file\label{ap:index.zeropoint}}

This index is used to store the result of the per-frame photometric zeropoint
determination. The format of the file is a header of the column names (with a hash
character in the first column) on the first line followed by a line entry for
each frame consisting of multiple colums with the form of \texttt{<FILENAME>
AIRMASS EXTINCTION FILTER MAG\_LIMIT NCALOBJS <ORACTIME> SKY\_VALUE
SKY\_VALUE\_ERROR SKY\_VALUE\_MAG TRANSPARENCY Z
EROPOINT ZEROPOINT\_ERROR ZEROPOINT\_SRC} where:
\begin{itemize}
\item \texttt{<FILENAME>} is the name of the reduced frame, 
\item \texttt{AIRMASS} is the mean airmass at the middle of the exposure, 
\item \texttt{EXTINCTION} is the extinction in magnitudes/airmass in the
specific filter (set to '-99' if no photometric statistics were determined), 
\item \texttt{FILTER} is the filter used,
\item \texttt{MAG\_LIMIT} is the limiting magnitude of the frame (set to '-99' if no photometric statistics were determined), 
\item \texttt{NCALOBJS} is the number of objects used in the zeropoint
determination (set to '-99' if no photometric statistics were determined), 
\item \texttt{<ORACTIME>} is the auto-inserted time of frame in standard ORAC
YYYYMMDD.ddddd format,
\item \texttt{SKY\_VALUE} is the sky value of the frame in ADU (set to '-99' if no photometric statistics were determined), 
\item \texttt{SKY\_VALUE\_ERROR} is the variance on sky value of the frame in ADU (set to '-99' if no photometric statistics were determined), 
\item \texttt{SKY\_VALUE\_MAG} is the sky value of the frame in magnitudes (set to '-99' if no photometric statistics were determined), 
\item \texttt{TRANSPARENCY} is the estimated transparency (0..1) of the frame
(set to '-99'' if no photometric statistics were determined),
\item \texttt{ZEROPOINT} is the zeropoint of the frame in magnitudes
(set to '-99'' if no photometric statistics were determined),
\item \texttt{ZEROPOINT\_ERROR} is the error on the zeropoint of the frame in magnitudes
(set to '-99' if no photometric statistics were determined),
\item \texttt{ZEROPOINT\_SRC} is the source of the zeropoint determination
(normally \texttt{<primitive name>-version}; set to 'N/A' if no photometric statistics were determined),
\end{itemize}

A section of an example file is shown below:
%\begin{sidewaystable}
\begin{terminalv}
#AIRMASS EXTINCTION FILTER MAG_LIMIT NCALOBJS ORACTIME SKY_VALUE SKY_VALUE_ERROR SKY_VALUE_MAG TRANSPARENCY ZEROPOINT ZEROPOINT_ERROR ZEROPOINT_SRC
lsc1m005-kb78-20150124-0351-e00_bp_ff 1.9601564 -99 gp -99 0 20150125.3628819 -99 -99 -99 -99 -99 -99 N/A
lsc1m005-kb78-20150124-0352-e00_bp_ff 1.90434185 -99 gp -99 0 20150125.3665509 -99 -99 -99 -99 -99 -99 N/A
lsc1m005-kb78-20150124-0364-e00_bp_ff 1.6940621 0.076346831 rp 15.87 7 20150125.3884028 49.1304748347854 10.8466273708444 17.2630570207312 0.316227766016838 24.65 0.168283194819187 _CALCULATE_ZEROPOINT_-r6338
\end{terminalv}
%\end{sidewaystable}

The index file is updated by filing new per-frame zeropoint determination results with the
calibration system in the \task{\_CALCULATE\_ZEROPOINT\_} primitive.

\end{document}
